\documentclass[11pt,a4paper]{article}
\usepackage[T1]{fontenc}
\usepackage[utf8]{inputenc}
\usepackage{lmodern}
\usepackage[margin=24mm,headheight=15pt]{geometry}
\usepackage{microtype,amsmath,amssymb,amsthm,mathtools}
\usepackage{booktabs,longtable,array,tabularx}
\usepackage{xcolor,enumitem,listings,fancyhdr,needspace}
\usepackage{titlesec,etoolbox}
\usepackage[hyphens]{url}
\usepackage[colorlinks=true,linkcolor=blue!40!black,citecolor=blue!40!black,urlcolor=blue!50!black]{hyperref}
\usepackage[nameinlink,noabbrev]{cleveref}
\definecolor{ink}{HTML}{17324D}
\definecolor{codeback}{HTML}{F4F6F8}
\titleformat{\section}{\Large\bfseries\color{ink}}{\thesection}{0.65em}{}
\titleformat{\subsection}{\large\bfseries\color{ink}}{\thesubsection}{0.65em}{}
\titleformat{\subsubsection}{\normalsize\bfseries}{\thesubsubsection}{0.65em}{}
\titlespacing*{\section}{0pt}{3.0ex plus 1ex}{1.3ex}
\titlespacing*{\subsection}{0pt}{2.4ex plus .6ex}{.9ex}
\setlength{\parskip}{.35em}
\setlength{\parindent}{1.2em}
\linespread{1.055}
\setlength{\emergencystretch}{2em}
\pretocmd{\section}{\Needspace{7\baselineskip}}{}{}
\setlist{itemsep=.18em,topsep=.4em,leftmargin=1.8em}
\pagestyle{fancy}
\fancyhf{}
\fancyhead[L]{\small\sffamily Leant2: answers to the expert questions}
\fancyhead[R]{\small\sffamily Technical investigation}
\fancyfoot[C]{\small\thepage}
\renewcommand{\headrulewidth}{.3pt}
\newcommand{\code}[1]{\texttt{\detokenize{#1}}}
\newcommand{\repo}{\textsc{Leant2}}
\newcommand{\Comp}{\operatorname{Comp}}
\newcommand{\foldr}{\operatorname{foldr}}
\newcommand{\cost}{\operatorname{cost}}
\newcommand{\FV}{\operatorname{FV}}
\newcommand{\Fin}{\operatorname{Fin}}
\newcommand{\Obs}{\operatorname{Obs}}
\newcommand{\type}{\ensuremath{\mathsf{Type}}}
\newcommand{\Prop}{\ensuremath{\mathsf{Prop}}}
\newcommand{\answer}{\paragraph{Answer.}}
\newcommand{\recommend}{\paragraph{Recommendation.}}
\newcommand{\status}[1]{\paragraph{Status.} #1}
\newtheorem{proposition}{Proposition}[section]
\newtheorem{theorem}[proposition]{Theorem}
\newtheorem{lemma}[proposition]{Lemma}
\theoremstyle{definition}
\newtheorem{definition}[proposition]{Definition}
\theoremstyle{remark}
\newtheorem{remark}[proposition]{Remark}
\newcolumntype{L}[1]{>{\raggedright\arraybackslash}p{#1}}
\lstset{basicstyle=\small\ttfamily,backgroundcolor=\color{codeback},
  frame=single,rulecolor=\color{black!15},framesep=5pt,
  breaklines=true,columns=fullflexible,keepspaces=true,
  showstringspaces=false,tabsize=2,aboveskip=.7em,belowskip=.7em}
\hypersetup{pdftitle={Leant2: Answers to the Expert Questions},
 pdfauthor={Technical research report prepared for Vladimir Reshetnikov},
 pdfsubject={Dependent program synthesis, Lean, search, carriers, evaluation and verification}}
\begin{document}
\begin{titlepage}
\thispagestyle{empty}
{\sffamily\color{ink}\large RESEARCH AND IMPLEMENTATION REPORT\par}
\vspace{18mm}
{\Huge\bfseries\color{ink} Leant2:\\[2mm] Answers to the\\[2mm] Expert Questions\par}
\vspace{8mm}
{\Large From further-improvement proposals to a sound,\\reproducible synthesis architecture\par}
\vspace{12mm}
{\large Prepared for Vladimir Reshetnikov\par}
{\large 21 September 2026\par}
\vspace{12mm}
\noindent\rule{\textwidth}{.6pt}
\vspace{3mm}
\noindent\textbf{Scope.} All 31 questions in research areas R1--R9 of
\path{docs/proposals/11-further-improvements}, together with their engineering
prerequisites, relevant Lean APIs, and the proposed acceptance gates.
\par\vspace{4mm}
\noindent\textbf{Repository snapshot.}\par
\noindent\small\path{784664ff34e819422510a4d470ab07fe48bb736b}\par
\noindent Lean toolchain: \code{leanprover/lean4:v4.34.0}.
\par\vspace{4mm}
\noindent\normalsize\textbf{Evidence boundary.} This is a source-based investigation,
with explicit mathematical arguments and eight executable finite-model checks.
It is not a claim that the proposed engine changes have been implemented or
that the repository's Lean benchmarks have been rerun.
\vfill
\noindent\small The accompanying archive contains this self-contained \LaTeX{}
source, its compiled PDF, a standard-library-only Python experiment, its results,
and build instructions. Recommendations are distinguished from established
results; negative answers state the exact bound or model to which they apply.
\end{titlepage}

\begin{abstract}
The further-improvements proposal identifies the right bottlenecks: interacting
metavariables, carrier invention, residual evaluation, indexed induction,
recursive-call synthesis, proof scheduling, retrieval, ranking, and optional
learned guidance. Several of its expert questions already have useful partial
answers in existing systems. Smyth removes trace completeness and supports
higher-order examples; Hoogle+ handles type classes through dictionaries;
Canonical demonstrates indexed synthesis and predicate synthesis with a supplied
generalized induction scheme; and Lean 4.34.0 already supplies a
proof-producing call-by-value evaluator. Other questions need a narrower
specification before a completeness or reproducibility claim is meaningful.

This report answers each of the 31 questions individually. Its central design is
a whole-state search with dependency-aware components, immutable ranking
objectives, typed residual observations, guarded caches, recursion capabilities,
and a single audited acceptance boundary. It gives a constructive finite-carrier
search, proves a sound observation-based state lower bound, explains the limits
of motive enumeration, and states sufficient conditions for conservative
pruning, independent-goal reuse, bounded synthesis completeness, and deterministic
replay. It corrects several consequential overclaims: open-goal counting is not
an admissible heuristic in general; incomplete observation rows are not states;
finite parametric tests do not cover arbitrary lengths; and sorting results
cannot undo deadline-dependent discovery. A staged implementation plan moves
existing Lean facilities and retrieval earlier, while keeping custom dependent
evaluation and unrestricted helper invention as separately testable research.
\end{abstract}
{\small\setlength{\parskip}{0pt}\tableofcontents}
\clearpage

\section{Executive conclusions and evidence}
\label{sec:overview}

\subsection{What should change first}
The highest-value next step is not to replace every subsystem at once. It is to
make the current engine's semantic boundaries explicit and then exploit
capabilities already present in its pinned Lean version. In particular, indexed
induction is not wholly absent from Lean's low-level API: the current
\repo{} search deliberately filters it out. Similarly, an equation-based,
proof-producing evaluator need not start as a new interpreter: \code{cbv} and
\code{decide_cbv} are present in Lean 4.34.0. These are concrete starting points,
not predictions of benchmark success \cite{repo,leaninduction,leancbv}.

The proposed architectural direction is:
\begin{quote}
Preserve a small checked acceptance boundary; make search states persistent and
self-describing; treat observations and proof obligations as typed constraints;
allow heuristics to change priority, but not the meaning of acceptance or of a
negative result.
\end{quote}
Three different guarantees must remain separate. \emph{Acceptance soundness}
means that a returned program has the requested type and checked contract under
the chosen axiom policy. \emph{Search conservativity} means that pruning and
memoization do not discard a permitted solution. \emph{Reproducibility} means
that a specified logical run can be repeated. A kernel check at the end protects
the first guarantee, not automatically the other two.

\subsection{Sources, versions, and what was actually checked}
The reviewed \repo{} snapshot is commit
\path{784664ff34e819422510a4d470ab07fe48bb736b}; its toolchain file specifies
Lean 4.34.0. The investigation read the proposal's TeX sections and inspected
\path{Leant2/Search/Core.lean}, \path{Leant2/Native/Transaction.lean},
\path{Leant2/Engine.lean}, \path{Leant2/Accept/Gate.lean}, and
\path{Leant2/Core/Types.lean}. Lean's induction and call-by-value tactic sources
were additionally checked at the \code{v4.34.0} tag
\cite{repo,leaninduction,leancbv}.

The proposal's performance evidence belongs to its own reported experiment at
commit \code{33cec8b}, not to a new run performed for this report. For example,
it reports complete scores on six historical acceptance harnesses, thirteen
unscored Church stretch cases that remain unsuccessful, and a particular profile
where residual evaluation accounted for 45\% of measured wall time. These are
valuable observations, but they do not establish that all future workloads are
evaluation-bound, or that a new cost model will preserve the same deadline-bound
candidate lists \cite{repo}.

No Lean executable was available in the report-generation environment. The
included Python checks were executed; no Lean patch, compiled Lean example,
benchmark speedup, model accuracy, or external-expert consultation is claimed.
Lean-shaped snippets below are design sketches unless explicitly identified as
quoting an inspected API. The mathematical propositions are proved in prose,
not machine-checked in this archive.

\subsection{Coverage of the questions}
Question identifiers below follow the order of the numbered expert questions in
each research area, rather than the order of the proposal's final expert table.
The question descriptions are paraphrases; their full answers have stable
R\emph{i}.\emph{j} headings.

\begin{center}\small
\begin{tabularx}{\textwidth}{@{}l c X@{}}
\toprule
Area & Questions & Principal outcome\\
\midrule
R1 Search & 4 & Dependency components; valid lower bounds; fixed objectives;
logical-work replay.\\
R2 Carriers & 5 & Bounded constructive search; sound state lower bounds;
parametricity conditions; fusion and motive proposals.\\
R3 Evaluation & 4 & Typed residuals; versioned demand caches; existing certified
CBV baseline; conditional angelic reasoning.\\
R4 Dependent types & 4 & Finite motive grammars; explicit transports; staged
transparency; existing induction APIs.\\
R5 Recursion & 3 & Capability realization theorem; top-down angelic search;
measurable scheme portfolio.\\
R6 Contracts & 4 & Proof-debt dependencies; a sufficient inductive fragment;
existing tactics; fair scheduling.\\
R7 Retrieval & 3 & Dependent hypergraph abstraction; dictionary passing already
in Hoogle+; leakage-controlled program benchmarks.\\
R8 Ranking & 2 & Declared preferences and MDL; certified equivalence versus
observation-based grouping.\\
R9 Learned proposals & 2 & A task-specific accuracy study; bounded, replayable,
non-authoritative proposal interface.\\
\bottomrule
\end{tabularx}
\end{center}

\section{Semantic boundaries and corrections}
\label{sec:foundations}

\subsection{The synthesis problem and its result types}
Fix an environment $E$, a local context $\Gamma$, a target $T$, a contract
$C:T\to\Prop$, a provider set $\mathcal L$, and a trust profile $\mathcal A$.
An accepted answer consists of closed abstractions of
\[
 \Gamma\vdash p:T,
 \qquad \Gamma\vdash \pi:C(p),
\]
with all transitive axiom dependencies allowed by $\mathcal A$. Where there is
no contract, the second component is trivial. An executable answer additionally
requires successful realization through the supported compilation path. That
last condition is not identical to kernel type correctness.

The inspected gate already performs synchronous kernel checking, rejects
unresolved expression holes and escaped locals, generalizes universe
metavariables jointly, and audits the transitive axiom closure. The engine's
acceptance callback constructs the contract proof type from the extracted
program, so the contract association is not merely left to an unrelated proof
argument. The standard profile permits \code{propext}, \code{Quot.sound}, and
\code{Classical.choice}; the strict constructive profile permits none of these
axioms \cite{repo}. Preserve this distinction when adding evaluators and proof
tactics: ``kernel checked'' does not mean ``axiom free.''

A result should distinguish at least the following statements:
\begin{enumerate}
\item A program and its contract proof were accepted.
\item A particular candidate, or every candidate actually visited, was refuted.
\item A precisely described finite grammar was exhausted.
\item A kernel-checked proof establishes impossibility for the actual target or
contract.
\item Work or time ran out without one of the stronger conclusions.
\end{enumerate}
The existing result algebra already recognizes most of these distinctions. In
particular, a count of rejected candidates must never become a universal
impossibility claim merely because all visited candidates failed.

\subsection{The invariant needed for safe pruning}
Let $s$ be a partial search state. Write $\Comp(s)$ for its permitted,
well-typed completions, including remaining universe constraints, local-context
requirements, and the declared grammar bounds.

\begin{definition}[Conservative refutation]
A pruning procedure is conservative if its return of \emph{refuted} implies
\begin{equation}\label{eq:refute}
  \neg\exists\theta\in\Comp(s),\ C(p_s\theta).
\end{equation}
Its return of \emph{unknown} asserts no impossibility.
\end{definition}

An evaluator that returns a plausible Boolean on an open expression is not
thereby a conservative refuter. It needs either a proof of the implication in
\eqref{eq:refute}, or a verified semantics whose soundness supplies that proof
schematically. Differential tests can discover violations, but agreement on a
large finite test set is not a proof of the implication.

\begin{proposition}[Safe use of conditional observations]\label{prop:conditional}
Suppose $A_s(\theta)$ is a set of inferred observations satisfying
$C(p_s\theta)\Rightarrow A_s(\theta)$ for every $\theta\in\Comp(s)$.
If the conjunction $A_s(\theta)\land C(p_s\theta)$ is impossible for all such
$\theta$, then pruning $s$ is conservative.
\end{proposition}
\begin{proof}
A completion satisfying the contract would satisfy the inferred observations
by the first hypothesis, contradicting the second. Notice that observations
need not hold for completions that violate the contract.
\end{proof}
This simple distinction is essential for recursive-call oracles: an example
specifying $f(a)=b$ constrains a recursive call of every \emph{valid} completion,
not of every syntactically legal completion.

\subsection{Corrections that affect the development plan}
\begin{longtable}{@{}L{.25\textwidth}L{.70\textwidth}@{}}
\toprule
Issue & Corrected statement and consequence\\
\midrule\endhead
Open-goal heuristic & A single unification can assign several holes. Summing a
positive cost over all open goals need not be a lower bound. Start with zero,
or with an explicitly justified relaxed problem.\\
Deadline determinism & Sorting accepted candidates fixes their order, not
which candidates were discovered. Use completed logical epochs and replay
receipts.\\
Carrier minimality & Minimize a declared representation-and-program cost, not
an unspecified ``smallest type.'' Do not allow the fold package to capture its
major input.\\
Observation rows & Missing entries are unknown. Distinct incomplete rows need
not denote distinct states. Use witnessed conflicts, plus transition
constraints.\\
Reverse as a fold & Reverse has a right-fold implementation with carrier
\code{List A}: append a singleton at each step. Failure of a small body grammar
is not nonexistence of that carrier.\\
Parametric tests & A generic position test can characterize one length under a
suitable parametricity theorem. Testing finitely many lengths does not
characterize all lengths.\\
Trace completeness & Smyth was designed to remove this requirement, and its
implementation supports higher-order examples \cite{smyth}.\\
Indexed synthesis & Canonical gives positive indexed and generalized-predicate
examples, though the generalized induction scheme is supplied in its example
\cite{canonical}.\\
Type-class retrieval & Hoogle+ already incorporates dictionary constraints and
instance components; higher-kinded variables are a stated limitation
\cite{hoogle}.\\
Existing evaluator & Lean 4.34.0 includes \code{cbv} and \code{decide_cbv}.
Benchmark these before building a replacement \cite{leancbv}.\\
Transport examples & For naturals, $n+0$ reduces to $n$; $0+n$ is the appropriate
simple example of an equality requiring propositional reasoning
\cite{leaninductbook}.\\
Regression gates & A changed ordering need not preserve deadline-truncated
candidate lists. Separate exhaustive bounded equivalence from deadline
performance.\\
\bottomrule
\end{longtable}

These corrections do not invalidate the proposal's overall direction. They
change which claims can be hard acceptance gates, which are experimental
hypotheses, and which facilities should be reused rather than reimplemented.

\section{R1: search with interacting metavariables}
\label{sec:r1}

\subsection{R1.1: alternatives to copying or sequentializing all goals}
\answer Use an AND/OR search over \emph{dependency components}, while retaining
whole-state snapshots as the correctness baseline. There is no general
representation trick that makes shared existential variables independent. The
useful alternative is to identify the portions that really are independent,
share immutable structure, and represent a component's answer as a guarded
substitution rather than as an unqualified ``goal solved'' flag.

Aesop supplies a practical best-first search precedent; its probabilities are
priority annotations, not guaranteed semantic costs \cite{aesop,aesopreadme}.
For \repo{}, the relevant unit is a store of constraints plus an obligation
list. The current implementation snapshots \code{Meta.SavedState} for an
alternative and commits only when the whole continuation succeeds. That is a
sound structural baseline to preserve during refactoring \cite{repo}.

\paragraph{What independence must include.}
Construct a graph whose vertices include obligations, expression metavariables,
universe metavariables, delayed assignments, local declarations, instance
obligations, and residual observations. Add dependency edges not just for holes
in the displayed target, but also for holes occurring in local types and values,
metavariable declaration contexts, delayed-assignment environments, and pending
contracts. An operation that may assign a shared metavariable induces a write
edge. Unknown effects conservatively connect components.

For example, a target with no visible metavariable can still be coupled to
another goal through a local instance whose type contains one. Conversely, two
goals may mention the same immutable constant without being coupled. Constant
sharing is a dependency on the frozen environment, not a shared writable hole.
Computing connected components after every node would be wasteful: maintain
incremental summaries, and recompute a component only when an assignment or
context extension touches its summary.

\begin{proposition}[A sufficient frame condition]\label{prop:frame}
Let two component solvers have read sets $R_1,R_2$ and write sets $W_1,W_2$ in a
fixed global environment. Suppose each solver is deterministic relative to the
store it reads; all relevant reads and writes are included; fresh names are
renamed apart; and
\[
 W_1\cap(R_2\cup W_2)=\varnothing,
 \qquad W_2\cap(R_1\cup W_1)=\varnothing.
\]
Then executing either solver does not change the other's applicability or
result, apart from fresh-name renaming, and their resulting substitutions can
be combined.
\end{proposition}
\begin{proof}
Every location read by the first solver has the same value before and after the
second solver, and conversely. Their writes are disjoint. Thus each run takes
the same steps and produces the same assignments as it would in isolation.
The union of the assignments is well-defined. The hypothesis that the read
sets include typing and scope dependencies makes the same argument apply to
the well-typedness of these assignments.
\end{proof}
The proposition is a semantic contract for an implementation, not a claim that
surface free-variable collection already satisfies it.

\paragraph{Memoization.}
A memo entry should contain a normalized goal-and-context fingerprint, the
relevant environment and transparency settings, a dependency fingerprint, a
solution expression, and any assignments it requires. On reuse, freshen private
metavariables, instantiate the context substitution, validate the guard, and
recheck the candidate. A memo holding one successful inhabitant does not justify
throwing away all other inhabitants when a later contract can distinguish them.
Cache a stream or family of alternatives, or prove that the component is
observationally irrelevant to all remaining obligations.

Initially memoize only closed, immutable-context subproblems. Then add guarded
open-context entries. This earns useful reuse before implementing a complete
dependency-aware AND/OR engine. Snapshot costs in the historical profile were
small; eliminating copying is not itself evidence that the total search becomes
faster \cite{repo}.

\subsection{R1.2: an admissible cost-to-close heuristic}
\answer There is no generally valid positive ``one unit per open goal'' rule.
Admissibility is relative to a fixed cost model. It requires
\[
 h(s)\leq\min_{p\in\Comp(s),\ C(p)}
       \text{remaining cost of producing }p.
\]
The cost of a search transition and the final ranking cost of a program are
also different quantities unless explicitly aligned.

A single application may unify a target $R\,?x\,?y$ with $R\,t_1\,t_2$ and assign
both argument holes while solving a proof hole. Charging separately for each
assigned hole can double-count the same transition. Exact-local steps can even
have zero cost in the proposed cost table. The supplied finite-model check
shows the basic obstruction with two obligations closed by one unit-cost action:
the sum heuristic is $2$ while the true residual cost is $1$.

The safest initial heuristic is $h=0$. A stronger one comes from an
\emph{optimistic abstract problem}: erase constraints, allow some assignments for
free, and solve the easier problem. To transfer its lower bound, every concrete
completion must induce an abstract completion of no greater cost. Abstract
hypergraph distances can then be used, but only after verifying this simulation
property for every search rule.

\paragraph{Combining bounds.}
For coupled goals, the maximum of individually valid lower bounds remains a
lower bound. Their sum need not. Summation is justified only for components
that require disjoint charged work, or by an explicit cost partition assigning
each concrete transition's cost among abstract subproblems without exceeding
that cost. Mere lack of shared metavariables is insufficient if a single
charged provider action is counted in several components.

For a true best-first minimum-cost claim, keep the incumbent's objective $J(p)$
and lower bounds $L(s)$ on all unfinished branches. Certification of optimality
requires
\[
 J(p)\leq\inf_{s\text{ on frontier}}L(s).
\]
Changing priorities with a learned score does not invalidate this criterion,
but stopping before it holds yields only ``best found.'' Zero-cost transitions
also need a finite closure or cycle control; otherwise a zero-cost layer can
prevent all positive-cost work from being reached.

\subsection{R1.3: just-in-time learning and ranking soundness}
\answer Learning from partial successes is a reasonable search heuristic, but it
is not a semantic justification for ranking. Probe provides an explicit
just-in-time learning precedent for bottom-up enumeration \cite{probe}.
Transferring its mechanism to dependent top-down search is an experiment: a rule
opens a state-dependent subtree, and the distribution of useful rules changes
as holes and examples are resolved.

Separate three fields in the design:
\[
  J(p)\quad\text{final objective},\qquad
  L(s)\quad\text{valid lower bound},\qquad
  P_e(s)\quad\text{priority in learning epoch }e.
\]
Only $P_e$ should change when the learner updates. Freeze $J$ for a query, or
version it visibly if the user changes preferences. Retain $L$ independently of
the learned predictions. A learner that prefers a subtree cannot establish that
the subtree contains the intended program, or that a disfavored subtree has no
better solution.

An implementable control is epoch-based adaptation: collect observations during
a fixed logical batch; aggregate them in a canonical order; update integer or
rational scores; freeze those scores for the next batch. Record the update data
and a model-version digest. Reserve a nonzero deterministic share for a complete
bounded baseline enumerator. This separates reproducible learning from
scheduler-dependent races and prevents permanent starvation.

The stopping receipt should say either ``best found under objective version
$v$'' or ``minimum certified over grammar $G_B$.'' A deadline-limited search is
allowed to return different valid programs under different priorities. What
would be misleading is to call those differences semantically optimal, or to
hide the changed objective from a ranking comparison.

\subsection{R1.4: wall-clock limits and machine-independent results}
\answer In general, a fixed positive wall-clock deadline, unrestricted machine
speeds, and a result that uses all work completed by that deadline cannot all
coexist with identical candidate sets.

\begin{proposition}[Deadline obstruction]\label{prop:deadline}
Suppose a deterministic run first discovers a new accepted candidate after
positive work $w$. On two machines executing the same logical run at rates
$v_1$ and $v_2$, a common deadline $D$ can fall between $w/v_2$ and $w/v_1$.
A policy returning all candidates discovered before the deadline then produces
different candidate sets.
\end{proposition}
\begin{proof}
Choose rates with $v_1D<w\leq v_2D$. Only the faster run reaches the discovery.
Sorting either set cannot insert the missing candidate in the slower set.
\end{proof}
A calibrated node budget is not a complete cure: different nodes may have very
different normalization, unification, or proof costs. Startup calibration also
makes the chosen logical budget machine-dependent.

Use two services instead. An \emph{interactive service} respects wall time and
reports its completed logical prefix. A \emph{replay service} runs the recorded
prefix, independently of how long it takes. Deterministic checkpoints can expose
only completed epochs; the number of completed epochs may differ by machine,
but a named epoch has a reproducible meaning. A fixed logical budget gives
identical results under deterministic execution when both machines finish it;
a hard earlier deadline can only return an explicitly incomplete run.

Parallel workers should commit completed batches in canonical order, rather
than cancel every competitor on the first result and hope that sorting fixes
the race. Use deterministic tie-breaks, stable provider order, fixed random
seeds, and recorded external proposals. General restart schedules can be useful
heuristics, but the classical Luby theorem for randomized Las Vegas restarts
is not automatically a theorem about a portfolio of correlated deterministic
Lean tactics \cite{luby}.

\section{R2: carrier discovery and fold synthesis}
\label{sec:r2}

\subsection{R2.1: can finite examples determine a minimal carrier?}
\answer Yes for a fully specified finite search problem; no unique semantic
``minimal Lean type'' follows from finite examples alone. The object to optimize
is a typed package, not just the spelling of its carrier:
\begin{equation}\label{eq:foldpackage}
 (K,z,s,q),\qquad z:K,\quad s:A\to K\to K,\quad q:K\to B,
 \qquad h(xs)=q(\foldr\ s\ z\ xs).
\end{equation}
Its cost should include the carrier, seed, step, finalizer, required helpers,
and proof obligations. Type-expression size, state cardinality, runtime storage,
and source-code size are different objectives. Isomorphic types can have quite
different library support and elaboration costs.

\paragraph{A scope restriction is essential.}
The package may depend on fixed parameters, such as a mapping function, but must
not capture the major input list $xs$ outside its fold. Otherwise take $K=1$,
ignore the fold, and let $q$ compute the whole answer from the captured $xs$.
Even without that capture, allowing an unrestricted finalizer merely moves the
original task into $q$: a list-valued carrier can store the entire input and an
arbitrary finalizer can process it later. That is a valid representation, but
not necessarily progress in synthesis. Bound and charge the entire package,
and disallow calls to the unfinished target except through approved recursion
capabilities.

\paragraph{A constructive finite procedure.}
Choose finite provider and literal sets, a finite carrier grammar at size bound
$b_K$, finite syntax bounds for $z,s,q$, and a fixed universe/context policy.
For a finite DSL with decidable typing and exact terminating evaluation, the
following procedure decides sample consistency at those bounds:
\begin{lstlisting}
for package in all_packages_in_nondecreasing_fixed_cost(bounds):
    if not well_typed(package): continue
    if every_sample_matches(package):
        return MinimumSampleConsistent(package, bounds)
return Exhausted(bounds)
\end{lstlisting}
Every syntactic candidate is visited; every check terminates; the first accepted
cost is minimal. These are the three facts needed for the guarantee. In a Lean
implementation, a bounded tactic or normalization attempt that times out must
produce \emph{unknown}, not a false typing or evaluation verdict. The engine
may then claim exhaustion only of the explicitly bounded operational search,
not of all semantically valid packages with the same type-size bound.

For finite alphabets and finite carriers, an even cleaner exact subproblem is
to enumerate total tables. With alphabet size $a$, state count $n$, and output
set size $b$, there are
\[
 n^{an}\cdot n\cdot b^n
\]
labelled choices of transition function, initial state, and finalizer. Evaluate
each table on the samples. An implementation may use SAT or a finite-domain
solver, but exhaustive enumeration already supplies a constructive decision
procedure. Reifying a satisfying table into a permitted Lean carrier is a
separate checked step; a table with $n$ states does not prove that an arbitrary
carrier grammar has an admissible representation of those states.

The included experiment exhausts all tables with one and two states for binary
parity on words of length at most four. It finds no one-state solution and two
labelled two-state solutions. A second experiment on unary length modulo three
finds none with one or two states and six labelled three-state solutions. These
are exact finite results, not evidence that the full carrier-invention problem
has become easy.

\paragraph{What finite examples cannot establish.}
A package consistent with the samples may fail on longer inputs. A minimal
sample-consistent machine need not be the minimal machine for the intended
function. Nor does the failure of one bounded body grammar prove that its
carrier is insufficient. ``No solution'' is meaningful only after specifying
which of these finite or universal problems was actually decided.

\subsection{A sound state lower bound from partial observations}
\label{sec:carrierbound}
There is a useful replacement for counting distinct incomplete observation
rows. Fix a finite set $V$ of suffixes. Join $x,y\in V$ by an edge when a prefix
$u$ has been observed with both $h(u\mathbin{++}x)$ and
$h(u\mathbin{++}y)$ defined in the sample table, and these outputs differ.
Call the resulting graph $G$ the \emph{witnessed incompatibility graph}.

\begin{theorem}[Observed state lower bound]\label{thm:statebound}
For a uniform fold package \eqref{eq:foldpackage} satisfying the observations,
the map $x\mapsto\foldr\ s\ z\ x$ assigns different carrier states to adjacent
vertices of $G$. Consequently, if $K$ has $n$ states, then
\[
 \chi(G)\leq n.
\]
In particular, a clique of size $r$ certifies that fewer than $r$ carrier states
cannot suffice for the observed behavior.
\end{theorem}
\begin{proof}
Suppose suffixes $x,y$ produce the same state $k$. For every prefix $u$, the
fold identity
\[
 \foldr\ s\ z\ (u\mathbin{++}x)
 =\foldr\ s\ (\foldr\ s\ z\ x)\ u
\]
shows that both concatenations produce $\foldr\ s\ k\ u$. Applying the same
finalizer gives equal outputs. Therefore a prefix witnessing unequal outputs
prevents equality of the two suffix states. The state assignment is a proper
coloring with at most $n$ colors.
\end{proof}
This proof uses only actual conflicting observations and fold composition. It
does not assume a completed observation table, a membership oracle, or a
particular grammar for the step.

For example, the three partial rows
\[
 (0,\ ?),\qquad (?,\ 0),\qquad (1,\ ?)
\]
are different as strings, but only the first and third are incompatible. Two
states can explain the table; counting three distinct rows as three required
states is unsound. Compatibility of partial rows is not transitive. Coloring
is a necessary condition, not sufficient for a transition-consistent machine:
transition congruence imposes additional constraints.

A cheap implementation can use a found clique rather than compute the exact
chromatic number. A hard prune against a proposed finite carrier should include
a checked cardinality bound and the conflicting sample witnesses. Counting
values \emph{constructible by a small expression grammar} is not a substitute
for counting all states the fold can reach, unless reachability inside that
bounded grammar is itself the stated model.

For comparison, full behavior defines the left-context equivalence
\[
 x\sim_h y\quad\Longleftrightarrow\quad
   \forall u,\ h(u\mathbin{++}x)=h(u\mathbin{++}y).
\]
Its quotient has well-defined prepend transitions and an output map. This is a
semantic minimal-state construction on reachable input states, but the quotient
can be infinite and need not have a decidable equality or a useful constructive
representation in the chosen grammar. Classical fold characterizations and
constructive refinements are directly relevant background: Weber and Caldwell
explicitly give evidence-dependent constructive transformations in Nuprl, not
an oracle extracting an implementation from arbitrary finite samples
\cite{gibbons,weber}.

\subsection{R2.2: polymorphic automata and the correct alphabet}
\answer Separate a genuinely parametric fragment from operations that inspect
or compare elements. For a pure relationally parametric function
$h:\forall A,\operatorname{List}(A)\to\operatorname{List}(A)$, the useful finite
observation at input length $n$ is its action on the list of distinct positions
$[0,\ldots,n-1]$ of type $\operatorname{List}(\Fin(n))$.

Assume the naturality property
\begin{equation}\label{eq:naturality}
 \operatorname{map}(f)(h_A(xs))
 =h_B(\operatorname{map}(f)(xs))
 \qquad(f:A\to B).
\end{equation}
Every length-$n$ list $xs$ is the image of the position list under the function
sending a position to its element. Thus the output at that generic position list
determines the output at every list of that length, including lists with repeated
elements. The argument also constrains whether the function duplicates, drops,
or reorders positions. It does not identify its behavior at every other length.

For any fixed $N$, the identity and the function that returns its input when
its length is at most $N$ and returns the empty list otherwise are both
parametric and agree on all tests up to length $N$. They disagree at $N+1$.
The finite-model experiment explicitly checks this at $N=4$. Therefore ``a
single monomorphic instance suffices'' must not be confused with ``finitely
many examples suffice.'' The polymorphic-testing literature addresses the
former issue under parametricity hypotheses \cite{polytesting}.

With an explicit equality dictionary, position identity alone is not enough to
model all observations: equality patterns of inputs matter, and the dictionary
must satisfy the laws used by the evaluator. For arbitrary Lean definitions,
polymorphic syntax alone is not a proof of \eqref{eq:naturality}. Classical type
inspection, unrestricted providers, or extra structure can invalidate the
intended free theorem. Require an audited syntactic fragment with a relational
interpretation, or a separately proved parametricity theorem for the components.

\recommend Start with generic-position observations as a \emph{proposal and
testing} mechanism for a restricted list DSL. Use finite equality-pattern or
register-style abstractions only when the supported equality operations are
explicit. Treat the extension to arbitrary dependent Lean types and instances
as a new semantics problem, not as a change of alphabet in an ordinary finite
automaton. Angluin-style exact learning also assumes access to specified query
oracles; a user's fixed sample set does not provide an equivalence oracle
\cite{angluin}.

\subsection{R2.3: backwards fusion and synthesis of auxiliary functions}
\answer Backwards fusion is useful when it generates \emph{checkable local
algebra equations}. It is not a general complete carrier-invention procedure.
Suppose $H=\foldr\ s\ z$ and seek $h=q\circ H$. For a known recursive
specification of $h$, try to solve equations relating $q(s(a,k))$ to the
specification's constructor case, together with $q(z)$ for the base case. A
successful solution must cover every reachable $k$, or a proved invariant that
contains all reachable states; finite examples alone do not establish this.

\paragraph{Reverse illustrates the distinction.}
A perfectly valid choice is
\[
 K=\operatorname{List}(A),\quad z=[],\quad
 s(a,k)=k\mathbin{++}[a],\quad q(k)=k.
\]
Induction on the input proves that the fold computes reversal. For the
endofunction representation one may instead take
\[
 K=\operatorname{List}(A)\to\operatorname{List}(A),\quad z=\mathrm{id},\quad
 s(a,k)=\lambda r.\,k(a::r),\quad q(k)=k([]).
\]
The stronger invariant $H(xs)(r)=\operatorname{reverse}(xs)\mathbin{++}r$
proves correctness. The first representation repeatedly appends to a growing
list and has quadratic list-traversal work under the usual linked-list cost
model. The second exposes accumulator-style computation. This is an efficiency
and grammar issue, not a proof that \code{List A} cannot be a carrier.

\paragraph{Useful completeness boundaries.}
For fixed finite algebras, fixed expression bounds, and exact local semantics,
exhaustively solving these equations is complete within the stated finite
space. For a list homomorphism, a fixed candidate binary operation and its
identity can be checked against the required algebraic laws. For polynomial
functors, a fixed carrier and finite collection of constructor operations lead
to analogous finite local obligations. Inventing the carrier, an invariant,
and arbitrarily complex operations is a different, unbounded problem.

A practical search should use failed equations as hints. When $q$ cannot recover
information needed by a constructor case, propose a product carrier containing
an additional statistic. When a result must be extended with an unknown tail,
propose a function-valued carrier. When two recurrences repeatedly recompute
related quantities, propose a tuple. Keep the former representation as an
alternative; a failed heuristic equation search is not a semantic refutation.
The constructive characterization literature is a source of proof obligations
and transformations, while the finite grammar is what gives the implementation
its actual decision boundary \cite{gibbons,weber}.

\subsection{R2.4: complete motive enumeration for primitive recursion}
\answer There is a simple complete \emph{schema} for a declared nondependent
primitive-recursive fragment with fixed parameter types. It is not a small
complete search procedure for all dependent primitive recursion. For
\[
 f:\mathbb N\to S_1\to\cdots\to S_k\to T,
\]
use the constant recursor result type
$K=S_1\to\cdots\to S_k\to T$. The successor branch receives an induction result
of type $K$, so recursive calls at the predecessor may use newly computed
parameter values. If result tuples or further accumulators are allowed, enlarge
$K$ by the corresponding product and function constructors.

Completeness for this fragment is conditional on enumerating all its base and
step expressions. With finite constructor/provider/literal sets and a size
bound, that is a finite grammar argument. Across increasing bounds, fair
enumeration is a semidecision procedure for expressible solutions, not a
terminating decision procedure for every failure. A cap such as carrier size
three is an engineering bound, not a completeness theorem for primitive
recursion.

Dependent motives $P(n,x)$ require a different grammar because the types of
parameters and results may change with $n$. Use a stratified sequence: constant
result, generalized fixed parameters, products, index abstraction, then
invariant-bearing dependent families. A finite target-occurrence grammar is
not closed under adding an accumulator: adding a new argument changes the
binding structure, not just an occurrence of an old index.

To control natural-number explosion, charge introduction of recursive structure,
not just the number of recursive calls exposed by a recursor. Use observed
numeric constants and a small literal basis; prefer already available arithmetic
providers; and bound nested eliminators separately from term size. Retain fair
later access to excluded numerals or schemes when claiming semicompleteness.
None of these controls establishes a universal polynomial-time search bound.

\subsection{R2.5: induction generalization, rippling, and lemma speculation}
\answer The transferable idea is to use \emph{failed proof alignment} as evidence
about a missing parameter or invariant. Do not transfer every heuristic as a
hard pruning rule. A branch where the induction hypothesis speaks about
$g(xs,a)$ but the recursive call requires $g(xs,a')$ suggests generalizing $a$.
A goal comparing a recursive result followed by an append with an unextended
hypothesis suggests an accumulator invariant. Repeatedly failing to relate two
computed statistics suggests tupling or an auxiliary lemma.

For \repo{}, implement a small proposal generator: compare a stuck branch goal
with available induction hypotheses, anti-unify their rigid structure, and
identify arguments whose differences are accounted for by a missing binder.
Compute the transitive dependency closure of any variable to generalize, then
attempt the new induction in an isolated snapshot. A proposed lemma receives
its own typed proof obligation and resource budget. It becomes reusable only
after kernel checking and axiom auditing.

Canonical offers an especially relevant partial precedent: its paper's
reverse-involution example supplies a generalized list-induction rule and the
solver synthesizes the needed predicate. That demonstrates that predicate
instantiation can be automated once the rule is available; it does not show
that the system automatically discovers every required rule
\cite{canonical}. This suggests a controlled first experiment: give both the
current engine and the improved engine the same small library of generalized
schemes, measure predicate discovery separately, then measure scheme proposal
as a second task. Otherwise failures of scheme selection, motive inference,
and body search are confounded.

\section{R3: residual and bidirectional evaluation}
\label{sec:r3}

\subsection{R3.1: higher-order and dependent live evaluation}
\answer Higher-order live evaluation is not an untouched starting point.
Smyth's core exposition uses first-order function examples for simplicity, but
its implementation supports higher-order function examples. Its hole closures
and resumption rules are directly relevant, as is the later Scrybe work on
example propagation \cite{smyth,scrybe}. The genuinely difficult extension for
\repo{} is dependent typing across changing hole assignments and contexts.

Represent a suspended occurrence as a \emph{typed closure}, not a bare hole
identifier:
\[
 (m,\Delta,\rho,T,\nu),
\]
where $m$ is declared in telescope $\Delta$, $\rho$ is a well-typed substitution
from that telescope into the evaluation context, $T$ is the expected result
family after substitution, and $\nu$ records the relevant state version. Two
occurrences of one hole under different environments are related constraints
on the same unknown term, not the same value.

For a function hole $g$, an observation can say that applying it to argument
$a$ must produce $b$. Multiple observations for $g$ must be simultaneously
realizable by one function. For a dependent function
$g:(x:A)\to B(x)$, the outputs live in different fibers $B(a)$. A table of
untyped values loses exactly the information needed to justify substitution,
comparison, and transport. One must preserve the expected fiber or use a
heterogeneous relation with explicit equality evidence.

\paragraph{The first implementation obstacle.}
A delayed metavariable assignment can mention a hole whose local telescope has
been changed by abstraction or recursion. The current
\code{instantiatePartial} explicitly documents that it can build expressions
whose remaining metavariable contexts no longer match, and restricts their use
to reduction rather than assignment \cite{repo}. This is not by itself a
demonstrated bug. It does mean that a proof of conservative residual evaluation
cannot simply appeal to a well-typed substitution lemma without first proving
that the partial-instantiation representation satisfies its premises.

Start with a typed first-order inductive fragment; add finite higher-order
application observations; only then add dependent families and transports.
When the supported fragment ends, retain a neutral expression and return
\emph{unknown}. Falling back must not silently reinterpret a dependent value as
a value in the wrong fiber.

\subsection{R3.2: incremental normalization and rollback invalidation}
\answer The right general precedent is dependency-tracked, self-adjusting
computation, not an assumption that a cache of normal forms remains valid
across rollback. Acar, Blume, and Donham give a correctness and consistency
semantics combining memoization and change propagation \cite{selfadjust}.
That supplies a conceptual model; it is not a ready-made Lean metavariable
normalizer with the required invariants already proved.

A residual evaluation node should record its expression or closure key, the
assignment and environment facts it actually read, its resulting value or
neutral residual, and edges to dependent nodes. On an assignment, invalidate or
resume only affected nodes. On rollback, restore the persistent graph version
or reject entries whose read fingerprints no longer match. Invalidating only
when a hole becomes assigned is insufficient: the same hole name can occur in
another branch with a different assignment, context, or delayed environment.

\begin{proposition}[Guarded evaluation reuse]\label{prop:cache}
Suppose an evaluator is deterministic and every semantic dependency of a cached
result is included in its guard. If the expression, context interpretation,
and every guarded dependency have the same meaning in a new state, reusing the
result is equivalent to reevaluating it there.
\end{proposition}
\begin{proof}
The evaluator sees the same expression and the same values for every inspected
input. Determinism therefore gives the same result. An implementation that
records too few dependencies does not meet the hypothesis.
\end{proof}

The guard includes more than assigned expression holes: universe assignments,
local declaration types and values, instance choices, equation-lemma sets,
transparency, options affecting reduction, and the environment fingerprint may
matter. Hashes are indices for these data, not logical equality proofs; verify
structural equality on collisions where correctness depends on it.

The current \code{SearchCtx} stores \code{residualBlockers} in an \code{IO.Ref}
outside \code{Meta.SavedState}; its checks examine whether remembered holes are
still present or assigned. That is an optimization to scrutinize when making
states first-class. A residual cache should be owned by its root program and
state version, whereas the monotone work ledger should deliberately remain
outside rollback. These two kinds of mutable state have different lifetimes
\cite{repo}.

The included rollback model demonstrates the failure mode, not a reproduced
\repo{} defect: caching value zero under one hole identifier and reusing it
after the same identifier denotes value one produces a stale answer. A
branch-and-assignment-aware key avoids that example. A production guard must
cover the complete dependencies, not merely mimic that toy key.

\subsection{R3.3: the cheapest conservative evaluator}
\answer First benchmark Lean's existing proof-producing call-by-value machinery.
The pinned version includes \code{Lean.Elab.Tactic.Cbv}, whose
\code{evalDecideCbv} applies \code{of_decide_eq_true} and calls the CBV goal
checker. The wrapper exposes \code{cbvGoal}, \code{cbvHyp}, and
\code{cbvDecideGoal} from the meta-level implementation \cite{leancbv}.
This makes a concrete fourth baseline alongside the current kernel reduction,
ordinary simplification, and a prospective custom evaluator.

The current reference manual describes CBV as equation-based propositional
rewriting, including support for well-founded definitions. It does not add
code-generator trust as \code{native_decide} does, but its standard-axiom use
still matters to a strict constructive profile. It also documents a limitation
on rewriting dependent positions \cite{leantactics}. Therefore it is a strong
candidate for closed or suitably nondependent observations, not a complete
solution to typed open-hole resumption.

\paragraph{A low-risk implementation ladder.}
First split contracts into independently reportable observations using proved
equivalences. Syntactic conjunction splitting is already implemented; extending
it to Boolean conjunction, a concrete list's \code{all}, and decidable equality
requires the corresponding reflection lemmas. In particular, an arbitrary
\code{BEq} instance does not justify replacing \code{x == y} by $x=y$ without
lawfulness evidence.

Second, apply ordinary proof-producing simplification or CBV to each observation,
with cached expression sharing and explicit resource limits. A failed or stuck
call is unknown. Third, introduce a small verified evaluator for the hottest
supported fragment. Its semantic theorem should state that every reported
concrete result agrees with the denotation of the quoted well-typed expression,
and every reported refutation excludes all permitted completions. Such a
fragment may exclude dependent eliminators while still improving many examples.

A fast untrusted interpreter can also propose a refutation certificate. Only a
successfully checked certificate causes a hard prune. If certificate generation
or checking fails, keep the branch. When proof production is too costly for every
node, the engineering choice is between a verified abstract transformer and a
heuristic signal that changes priority only. There is no third option in which
an unverified false result becomes safe merely because surviving candidates
later pass the gate.

\paragraph{What Lean4Lean contributes.}
The current Lean4Lean paper reports a complete alternative typechecker and
correctness proofs for parts of the kernel, not a fully verified, automatically
extractable open-hole evaluator for arbitrary dependent synthesis
\cite{lean4lean}. Reusing its definitions or proved lemmas can reduce duplicated
metatheory. Porting rules alone does not transfer a theorem about different
closure representations, partial substitutions, or the host evaluator. State
the simulation relation first; then identify which existing lemmas actually
prove its cases.

\subsection{R3.4: top-down angelic search and completeness}
\answer Top-down angelic reasoning is feasible, but Burst's theorem does not
transfer automatically. Burst's completeness statement assumes a complete
angelic synthesizer and valid explanations of failure. Its outer algorithm also
checks candidates under ordinary semantics and backtracks over strengthenings
\cite{burst}. A top-down implementation must establish the same kinds of
conditions for its own grammar and constraint representation.

For a recursive call at a known argument, an angelic semantics permits an output
consistent with the current contract information. Equal calls of one function
must agree when their arguments are known equal; dependent calls must return
values in the correct fibers. Ignoring these restrictions is an overapproximation
and may merely waste search. Enforcing an unjustified equality or restricting an
unknown output too narrowly can instead remove a real solution.

Use the following contract:
\begin{quote}
Every ordinary execution of every valid completion has a corresponding angelic
execution admitted by the residual constraint store.
\end{quote}
Unsatisfiability of that store is then a conservative pruning reason. A satisfying
angelic store is only permission to continue: its separately chosen outputs may
not arise from a single actual recursive program. Reevaluate the completed body
under ordinary semantics and subsequently pass the full acceptance gate.

\begin{proposition}[Bounded top-down completeness]\label{prop:bounded}
Fix a finite typed candidate grammar, exact terminating sample evaluation, and
a search whose branch pruning is conservative. Suppose the search enumerates a finite set of bounded derivations,
visits every unpruned candidate, terminates after that enumeration, and no
refinement permanently discards a sample-correct candidate. Then search terminates and finds a sample-correct candidate whenever
one exists in that grammar; otherwise it exhausts the grammar.
\end{proposition}
\begin{proof}
There are finitely many candidates. Each evaluation terminates. Conservativity
and refinement preservation leave every sample-correct candidate available,
and fairness eventually visits it. If none exists, exhausting the finite set
terminates with failure at the declared bounds.
\end{proof}
The assumptions contain the real work. An unbounded value domain, an infinite
refinement sequence, or a fuel-limited evaluation returning unknown does not
satisfy the exact finite premises. Fuel exhaustion is not a refutation. A
practical bounded system can ensure progress by blocking an already disproved
candidate, retaining unresolved alternatives, and recording which operational
bounds prevented a complete verdict.

\section{R4. Indexed families and motive synthesis}
\label{sec:r4}

\subsection{R4.1. Can motive enumeration be complete, terminating, and closed under generalization?}
\answer Yes for a precisely bounded syntactic class, but the class of occurrence
abstractions alone is not closed under accumulator invention. Three different
operations should be separated: abstracting indices already in a target;
reverting existing parameters into the induction hypothesis; and inventing a
new carrier or invariant not already present in the context. The first is
largely elaborator engineering. The second is a small but important search.
The third contains the genuine synthesis problem.

For example, mapping a length-indexed vector needs the motive
\[
 M(n,v)=\operatorname{Vec}\,B\,n.
\]
An accumulating traversal may instead need
\[
 M(n,v)=\prod_{a:S(n,v)}T(n,v,a),
\]
or an invariant-carrying result. Selecting some occurrences of $n$ in the first
expression cannot invent the new binder $a$ or its type. Consequently, even a
complete search through maximal-occurrence abstractions is incomplete for the
second class.

A useful finite specification is a \emph{motive template language}. Its terminals
are scoped subexpressions of the target, local types, and an explicitly bounded
set of type constructors. Its operations are dependency-closed abstraction,
reversion of a bounded number of locals, and bounded products, arrows, and
subtypes. Bound expression size, universe-expression size, and the number of
new binders. Enumerate the resulting finite syntax, discard ill-scoped terms,
and check each remaining candidate in the appropriate telescope. If checking
and the permitted constraint solving terminate exactly on this fragment, the
procedure is complete and terminating for that grammar. Operationally timed-out
checks must remain unknown rather than becoming evidence of nonexistence.

Dependency closure is essential. Abstracting an index changes the types of
later variables that mention it. Those variables must be transported or
reabstracted as well. One cannot independently replace two textually identical
occurrences when their binding or type dependencies differ. Nor should
``maximal occurrence'' mean ``largest printed substring'': the operation is on
scoped expressions and their dependency graph.

\recommend Start with one inferred motive for a local whose indices are distinct
variables, then a short ordered list of explicit generalizations. Keep genuinely
invented motives in the same bounded carrier/invariant service used by R2 and
R6. This gives a clear progression from engineering to research and avoids
charging a basic vector-map example for an unrestricted higher-order search.
The existing Canonical examples also show that indexed synthesis is not an
entirely unexplored capability, although they do not establish completeness for
arbitrary motives \cite{canonical}.

\subsection{R4.2. How should transports be represented and canonicalized?}
\answer Use typed transport nodes with checked equality evidence; normalize a
selected index language, not all propositional equality. There is no general
cheap canonical representative modulo arbitrary propositions that Lean may
prove. Replacing this problem by heterogeneous equality or a setoid changes its
presentation, not its underlying proof obligations.

If $P:I\to\type$, $x:P(i)$, and $e:i=j$, the transport has type $P(j)$:
\[
 \operatorname{cast}(\operatorname{congrArg}P\ e,x):P(j).
\]
The indices $i,j$ and family $P$ must remain in the internal representation.
Proof irrelevance can identify two proofs of the \emph{same} equality proposition;
it does not identify arbitrary endpoints or make all transports disappear
computationally. A display layer may hide a checked cast, but the term submitted
to the kernel must still contain the required transport unless conversion makes
it unnecessary.

A practical normalization discipline is to select a terminating, deterministic
rewrite system for an index fragment and produce an equality certificate with
every rewrite. Normalize each endpoint using that policy, compose the
certificates, and reuse one transport construction for that normalized pair.
Unknown index equations remain proof obligations. This yields canonicalization
\emph{relative to a policy}, not modulo all provable equalities. The policy and
its lemma set belong in cache keys and receipts.

The proposal's example should be reversed: for Lean's natural-number addition,
$n+0$ is definitionally equal to $n$, whereas $0+n=n$ generally needs a proof
when $n$ is a variable \cite{leaninductbook}. Therefore
\code{Vec A (n + 0)} is not a good test of propositional transport. Use, for
example, a vector at index $0+n$ or an associative rearrangement of an index
expression.

Heterogeneous equality is useful for stating that values in different fibers
correspond. It is not a universal congruence engine. Even the current Lean
\code{cbv} documentation explicitly identifies restrictions on rewriting in
dependent positions; changing to heterogeneous equality does not automatically
supply the missing dependent congruence principles \cite{leantactics}.
For the first implementation, restrict transport to a registered family and a
small index theory. Charge proof and normalization work separately from the
presentation-size score, rather than calling every cast literally free.

\subsection{R4.3. What transparency policy is supported by existing experience?}
\answer The inspected systems support selective normalization, not indiscriminate
full unfolding. Canonical describes unfolding aliases of function types,
reducible definitions, and type-class instances in its translation. The
\code{sauto} interface distinguishes heuristic unfolding from explicitly forced
unfolding, with special treatment of basic logical connectives
\cite{canonical,sauto}. Neither description establishes one globally optimal
transparency policy for \repo{}.

Use three explicit levels. At the first level, compute weak-head forms needed to
recognize binders, inductive heads, and obvious reducible aliases. At the second,
unfold a specific definition whose head is blocking an application, an index
comparison, or motive construction. At the third, permit a larger but bounded
normalization attempt, with a separate work charge and cache namespace.

A failed first-level comparison is not proof that two types differ. It is a
request either to postpone a constraint or to try the next permitted level.
Conversely, normalizing every definition before retrieval can erase useful
library boundaries, duplicate large expressions, and spend work on arguments
whose values are irrelevant to the current goal. Preserve high-level function
names for retrieval and presentation even when a checked reduced form is cached
for conversion.

The simplest experiment compares fixed reducible-head normalization against this
three-level policy on identical logical search traces. Record the number and
cost of successful escalations, expression growth, cache hit rate, and candidates
lost only to resource limits. Do not choose a policy from one profile whose
runtime was dominated by residual contracts rather than conversion.

\subsection{R4.4. Can Lean's motive computation be called as a library?}
\answer Yes. The pinned Lean 4.34.0 source exposes relevant operations in
\path{Lean/Meta/Tactic/Induction.lean} \cite{leaninduction}. In particular:
\begin{description}
\item[\code{Lean.Meta.getMajorTypeIndices}] retrieves indices subject to the
low-level induction routine's restrictions, including variable-form indices and
checks on repeated or improperly dependent occurrences.
\item[\code{Lean.Meta.mkRecursorAppPrefix}] constructs a recursor application
through its motive, abstracting the major premise and indices from the target.
\item[\code{Lean.MVarId.induction}] orchestrates reversion, reintroduction,
motive construction, and creation of the constructor subgoals. Returned
\code{InductionSubgoal} records include the new goal, fields, and a local-variable
substitution.
\end{description}
These operations compute a motive for the prepared target; they do not discover
an arbitrary stronger invariant. Their availability answers the library-access
question, but not the empirical question of whether repeated calls are cheap
enough on a particular workload.

The current \repo{} rule calls \code{MVarId.induction} after requiring zero
indices and excluding \code{Nat}. Those exclusions are in the search rule, not
a statement that Lean lacks the corresponding induction facilities
\cite{repo}. A small first experiment should admit an indexed family only when
its indices satisfy the low-level preconditions and its recursor is supported.
For a non-variable index such as $n+1$, use a prepared generalization with a
remembered equality or an elaborator-level induction path that handles the
expression; do not pass an incompatible major premise to the low-level routine
and interpret the resulting exception as logical impossibility.

Wrap every attempted preparation in a transaction. Classify failure as
unsupported index shape, prohibited elimination from \Prop, failed type
constraint, unsupported recursor, or exhausted resources. Preserve these
categories in traces. Cheap syntactic prechecks can avoid predictable failures,
but no particular millisecond bound is established by this report. The wrapper
should also expose the resulting dependency substitution to the search rather
than reconstructing it from printed variable names.

\section{R5. Recursive programs from incomplete examples}
\label{sec:r5}

\subsection{R5.1. Does the capability interpreter agree definitionally with its realization?}
\answer Not in general. A structural realization can have useful definitional
computation rules, but well-founded recursive equations are often available only
as propositional equalities. Lean's documentation explicitly distinguishes the
two cases \cite{leanrec}. The correct interface should promise a proved
simulation property, not universal definitional equality.

For a recursive state $x:S$, a relation $R:S\to S\to\Prop$, and result family
$P:S\to\type$, the fully dependent capability is
\[
 \mathit{rec}_x : \prod_{y:S} R(y,x)\to P(y).
\]
Include every changing parameter in the recursive state or in the family over
which the induction hypothesis is generalized. A capability that mentions only
the original major argument is insufficient when recursive calls change an
accumulator, an additional list, or a dependent result index.

Suppose a typed body language has an interpreter
\[
 \operatorname{eval}(k,b,x):\operatorname{Option}(P(x)),
\]
where $k$ is fuel, and a realization $\operatorname{realize}(b)$ obtained from a
checked well-founded recursion construction. The desired one-sided theorem is
\begin{equation}
 \operatorname{eval}(k,b,x)=\operatorname{some}(v)
 \quad\Longrightarrow\quad
 \operatorname{realize}(b)(x)=v.
 \label{eq:simulation}
\end{equation}
Its assumptions must include well-typed bodies, valid decreasing-call
certificates, sound interpretation of every primitive, and the equations used
by the realization. Prove it by induction on the successful evaluation or by
well-founded induction, using the realization's equation theorem at each call.
The interpreter's \code{none} is an inconclusive resource result, not the value
of the realized function and not evidence that a contract is false.

This interface is deliberately weaker, and more useful, than claiming the
interpreter computes exactly what \code{Kernel.whnf} does on a raw
\code{WellFounded.fix} term. Once an interpreted observation is checked against
\eqref{eq:simulation}, it can support a theorem about the realized function even
when conversion alone cannot perform the same computation.

\recommend Benchmark existing equation-based \code{cbv}/\code{decide_cbv}
before implementing a new interpreter. Keep direct compiled evaluation as a
screening or ranking facility unless its result is independently certified.
In particular, \code{native_decide} has a different trust story from ordinary
kernel reduction: its proof uses a compiler-related axiom, so it is not
interchangeable with the current standard axiom profile
\cite{leancbv,leantactics,repo}.

After realizing a candidate, check its actual contract again and separately
check that the intended executable definition compiles. A theorem about one
kernel term is not automatically a correctness theorem for an arbitrary
\code{implemented_by} replacement. An equality or simulation linking the
published implementation to the certified term must be checked and recorded.
If compilation fails, preserve the verified logical object but report the
execution failure rather than silently changing the implementation.

\subsection{R5.2. Does angelic execution compose with top-down search?}
\answer Yes as an over-approximate constraint semantics; the challenge is retaining
all relevant alternatives efficiently. Top-down search contributes expected
types, lexical contexts, and dependent constraints early. A bottom-up pool
contributes already closed expressions, reusable behavior vectors, and a
convenient finite representation for combinations. Moving from one organization
to the other does not transfer a completeness or termination proof for free.

The appropriate top-down object is a typed hole together with a set of possible
behaviors in the environments where it is observed. A recursive-call table must
respect repeated calls to the same input and the result fiber for that input.
An arbitrary independent answer at every occurrence is usually too weak; one
fixed guessed answer with irreversible commitment is too strong. Maintain the
constraint relating those answers, and backtrack the assumptions with the
program branch.

The preservation condition and finite completeness statement of
\cref{prop:bounded} apply here. An angelic success only schedules further work.
A closed candidate is evaluated ordinarily. If it contradicts an assumed call
result, reject that candidate or refine its assumption store without excluding
other bodies that could satisfy the contract. If the call's result was already
a consequence of the contract, the stronger conditional refutation argument in
\cref{prop:conditional} is available. These two sources of information must have
different tags.

A small hybrid is preferable to replacing the full search. Enumerate shallow,
closed terms bottom-up for a specific local telescope and expected type, then
offer them as top-down alternatives. Reuse across contexts requires explicit
renaming and a checked substitution. A behavior vector is an index into this
pool, not permission to merge terms forever: later recursive calls can observe
an input on which previously indistinguishable terms differ.

This recommendation uses the mechanisms discussed in R3 rather than assuming
that the existing Burst implementation already proves the desired top-down
algorithm correct. The gains or losses in memory, propagation power, and search
coverage require the same typed benchmark and work budget on both variants.

\subsection{R5.3. Is there a canonical small basis of recursion schemes with measured coverage?}
\answer There is no established canonical optimum for the combined Lean corpus
specified here. A claim about coverage requires a fixed provider library,
representation, cost bound, and meaning of ``the same problem.'' Existing
recursive-synthesis corpora and systems are useful sources of tasks, but their
success rates cannot be assigned to a new scheme basis without rerunning it.
Trio is a relevant example of constructing recursive programs from nonrecursive
expressions rather than merely enumerating raw recursor terms \cite{trio}.

A defensible initial basis is the following. The examples describe expressibility,
not newly measured synthesis success.
\begin{longtable}{@{}L{33mm}L{56mm}L{59mm}@{}}
\toprule
Scheme & Typical witness & Additional obligation\\
\midrule\endhead
Single structural fold & List map, length; tree size & Constructor algebra and
base value\\
Generalized parameter or function carrier & Accumulating reverse; left-fold
behavior expressed by a right fold & Generalized induction hypothesis and
finishing function\\
Tupled structural result & Fibonacci using a pair; count and sum together &
Product algebra and projection\\
Multiple-argument recursion & Zip; merge with a decreasing total size &
Generalization or a relation on the joint argument state\\
Course-of-values recursion & Calls to deeper, certified predecessors &
Accessible predecessor capability and correct indexing\\
Well-founded measure & Euclidean-style recursion and decreasing search intervals &
Well-foundedness and a decrease proof at each call\\
\bottomrule
\end{longtable}
The rows overlap intentionally. Fibonacci does not intrinsically require
course-of-values recursion: primitive recursion into a pair suffices. Zip can
recurse on one list while inspecting the other. A flattening function may become
easy once append is a permitted component, without inventing a new recursion
scheme. These distinctions prevent attributing every failed body search to a
missing recursor.

For each task record the cheapest sufficient scheme under a declared provider
set, initially using a hand-written witness. Compare a carrier-given track, a
scheme-given track, and a fully inferred track. The difference between their
scores separates body-search difficulty from scheme invention. A proposed basis
is successful if its measured coverage and cost improve; it need not be called
canonical. Keep an explicit fallback to a larger bounded grammar so that a
convenient initial basis does not become an undocumented completeness claim.

\section{R6. Contracts, induction, and proof debt}
\label{sec:r6}

\subsection{R6.1. How should value holes and proof holes be interleaved?}
\answer Schedule a dependency graph of obligations, not two global queues called
``program'' and ``proof.'' A proof about unresolved code is sometimes useless,
but sometimes it directly determines that code. Reflexivity can assign an index
hole; an equality constraint can rule out a constructor before its arguments
are filled; a type-class obligation can determine a type parameter. Therefore
the blanket rule ``proofs wait until all values are closed'' loses important
information.

Partition obligations by what they can do now. Eager obligations include
well-scoping, type conversion, index constraints, required instances, and cheap
local contradictions. Once their relevant value dependencies close, local
branch contracts become eligible for a small proof service. Expensive global
properties, induction choices, and speculative lemmas accrue proof debt and
wait for a promising value skeleton. Every debt record carries its owning
candidate, local telescope, dependencies, already tried proof tiers, and the
snapshot or substitution under which it was created.

Synquid establishes the value of decomposing refinement specifications into
component obligations \cite{synquid}. The transferable idea is decomposition,
not a requirement to put every Lean proposition into an external SMT solver.
For \repo{}, start with a small collection of checked abstract facts: length
expressions, numeric bounds, constructor information, and selected membership
facts. Each transfer rule has a Lean theorem. A local abstract calculation then
proposes a proof composition that the kernel checks. Unknown abstract facts
stay unknown.

For example, synthesize vector-like list mapping together with length
preservation using a motive of the form
\[
 M(xs)=\{ys:\operatorname{List}B\mid |ys|=|xs|\}.
\]
In the cons branch, the induction hypothesis supplies a recursive value $ys$
and a proof $|ys|=|xs|$. Constructing $f(a)::ys$ leaves only
$1+|ys|=1+|xs|$, which follows locally. There is no need to wait until a large
closed recursor term is assembled before discovering the same induction.

This route is not a universal replacement for global verification. Involution,
commutation of two calls, or preservation under concatenation may require a
stronger invariant or an auxiliary lemma. The proof-debt graph should expose
that difference rather than repeatedly spending the same local tactic budget.

A concrete integration issue is already visible: the current
\code{proofPortfolio} rejects a target containing free variables, in addition to
metavariables \cite{repo}. Thus the proposed local-branch proof service cannot be
obtained merely by adding another tactic name to its existing list. It needs a
context-aware wrapper, or a closed abstraction over the local telescope, with
correct assignment and replay of the resulting proof. Test this explicitly on
an induction hypothesis used inside a branch.

\subsection{R6.2. Which properties admit induction, simplification, and arithmetic?}
\answer A useful sufficient class can be recognized, but there is no reason to
expect an exact syntactic classifier for every property that some future
induction-and-rewriting search might prove. Define the class by a certificate
format rather than by an informal expectation of simplicity.

Fix a structural recursion scheme, an invariant template, a terminating rewrite
policy, and a declared decidable theory for the remaining side conditions. A
verification-condition generator creates one obligation per constructor. The
certificate records the induction scheme, the invariant instance, the rewrite
proofs, and the decision-procedure certificates. A successful translation into
this format establishes membership in the supported class. Failure to translate
means unsupported, not false.

\begin{proposition}[Local verification conditions suffice]
\label{prop:localvc}
Let $I$ be an inductively generated input type with a sound induction principle,
and let $f$ be defined by its structural recursion. Suppose that, for every
constructor, the corresponding branch satisfies $P$ assuming $P$ for precisely
the recursive arguments supplied by that induction principle. Then
$\forall x:I,\ P(x,f(x))$.
\end{proposition}
\begin{proof}
Apply the induction principle for $I$. At each constructor, unfold the defining
equation for $f$ and apply the certified branch obligation to the induction
hypotheses. These are exactly the premises required by that obligation.
\end{proof}
The substantive automation task is generating and proving the branch obligations,
not the final induction theorem. Length preservation of map falls directly into
this format. Reverse involution may need an append or accumulator invariant;
sorting generally requires order and permutation lemmas. Such examples should
be labeled by their additional certificate requirements rather than counted as
unexpected failures of arithmetic.

Presburger arithmetic is a sensible index theory, but distinguish the theory
from a particular tactic's implementation and timeout. The current Lean manual
states that \code{omega} does not yet implement every component of a full
decision procedure, while describing a complete linear-integer-arithmetic
procedure inside uncapped \code{grind}; capped invocations can still fail
\cite{leantactics}. Therefore use the phrase ``the registered, certified
arithmetic fragment'' in a completeness claim unless all required solver
premises have actually been established. A timeout is an unresolved obligation.

The same idea supports a bounded invariant search: enumerate templates, generate
finite branch conditions, and run exact fragment checkers when available. This
is complete only for the template language and solver fragment, not for all
true properties of structurally recursive programs.

\subsection{R6.3. Which Lean induction facilities are usable now?}
\answer Assemble existing facilities around an explicit induction plan. Lean's
low-level induction routine is available in the pinned version, and the current
reference also documents \code{fun_induction} for eligible structurally or
well-founded recursive definitions, using function-specific induction theorems
\cite{leaninduction,leantactics}. These facilities reduce the cost of building a
proof service, but do not themselves invent every required generalization or
lemma.

The recommended adapter has three stages. First prepare the induction variable
or recursive equation and a small ordered set of generalizations. Next obtain
constructor or function-induction subgoals. Finally simplify computation rules
and dispatch local obligations to an axiom-audited portfolio. Use direct
assumptions, reflexivity, selected simplification, and decidable evaluation
before larger arithmetic or saturation attempts. Library lookup belongs before
repeated speculative induction when a known theorem already solves the goal.

Aesop supplies configurable proof search and normalization; \code{grind}
combines its documented saturation and theory-solving facilities; premise
selection and hammer tools can supply useful lemmas
\cite{aesopreadme,leantactics,leanhammer}. None of those facts supports a promise
that arbitrary induction with helper-lemma invention is now solved in Lean.
In particular, adding \code{grind} to the existing closed-goal portfolio does not
replace choosing an induction hypothesis strong enough for the program.

For the implementation, keep version-specific tactic calls behind adapters.
The exact presence and behavior of every optional tactic must be tested against
the project lockfile, rather than inferred from a moving online manual. The
pinned induction and \code{cbv} entry points were source-verified here; no full
integration build was performed. This separation permits a core-only default
and optional external proof services without changing the acceptance gate.

An experiment on twenty or more quantified contracts should retain hand proofs
classified as: direct simplification, arithmetic, structural induction,
generalized induction, library lemma, or invented lemma. Report success and cost
by class. An aggregate solved count conceals whether a new service actually
addresses the intended induction bottleneck.

\subsection{R6.4. Is there a principled policy for scheduling proof debt?}
\answer There is a principled optimization problem, but no justified claim that a
standard independent-arm bandit theorem solves this particular scheduler.
Proof attempts share lemmas, metavariables, counterexamples, and partially solved
obligations. Their rewards change when a value choice changes. They also have
switching costs from context setup and normalization. Treat expected
benefit-per-cost as a heuristic, and keep fairness and acceptance semantics
independent of its estimates.

For an eligible obligation $o$ and proof tier $t$, a useful scheduling signal is
\[
 U(o,t)=\frac{\widehat{p}_{o,t}\,V(o)+\lambda\,I(o)}
                  {\widehat{c}_{o,t}+s(o,t)}.
\]
Here $\widehat{p}$ estimates success probability, $V$ estimates how much candidate
work a proof unlocks, $I$ estimates information useful for rejecting or refining
a branch, $\widehat{c}$ is work cost, and $s$ is switching overhead. This formula
is a proposed policy, not an optimality theorem. In particular, learning its
parameters must not turn a failed prediction into a refutation.

A robust first scheduler uses deterministic rounds. Give every eligible
obligation a cheap tier, select promising survivors for a larger tier, and age
unselected obligations so that a resumable unbounded run does not starve them.
Batch attempts that share a local context when the saved setup cost is material.
Use logical work quanta, not a fixed millisecond ladder, for replay; expose wall
clock as a separate interactive cutoff. Record attempted tiers so rollback does
not repeat an identical expensive failure without a changed dependency.

Counterexamples are particularly valuable shared information. For a universally
quantified contract, a closed, well-typed counterexample satisfying all
preconditions can refute a candidate and become an early test for later ones.
A malformed generated input, a failed premise, or an evaluator timeout cannot.
The candidate still needs a proof of the whole universal contract before
acceptance. This preserves the useful CEGIS discipline without confusing
repeated testing with certification.

\section{R7. Component retrieval at library scale}
\label{sec:r7}

\subsection{R7.1. What abstraction works for dependent, universe-polymorphic signatures?}
\answer Use a layered, constraint-carrying abstraction of \emph{applications}, not
an ordinary graph that connects each argument type separately to the result.
An application requires all of its arguments. A provider $A\to B\to C$ does not
make $C$ constructible merely because an $A$ is available. The appropriate first
approximation is an AND-hypergraph or a resource model with conjunctive
prerequisites. The accompanying finite model includes this elementary but
consequential counterexample.

An abstract signature should retain the sort, head constant, binder arity,
explicit/implicit/class roles, and sharing relationships between type variables.
At a finer level it retains selected rigid indices and simple universe
constraints. For example, the signature
\[
 (n:\mathbb N)\to \operatorname{Vec}\,A\,n\to B
\]
requires the vector's index to agree with the chosen first argument. Keeping
only the marginal heads \code{Nat}, \code{Vec}, and $B$ deliberately forgets this
relationship. Such forgetting may create a false positive, which concrete
elaboration must reject; it must not silently justify an executable application.

Define sound abstraction relative to a fixed candidate language: every concrete
application derivation allowed by that language must map to an abstract
derivation. This implies that failure of exact abstract reachability excludes
a concrete derivation, provided the abstraction and search bound correspond.
Success remains only a proposal. Higher-order introductions, locally available
functions, partial application, and dictionary construction must also be
represented if they are allowed in the concrete language. An abstraction that
models only first-order calls is not a sound exclusion oracle for a higher-order
search.

Use refinement to restore information precisely where a candidate path failed:
a shared index, a constructor of an indexed term, a universe relation, or a
class prerequisite. Refuting one instantiation does not refute a provider for
every instantiation. Keep refinement constraints scoped to the abstract
combination or query that justified them. The general TYGAR architecture gives
a well-motivated starting point, but the dependent refinement relation itself
requires a new correctness argument \cite{hoogle}.

Finally, distinguish \emph{retrieval ranking} from \emph{sound exclusion}. Taking
the top twenty components is a heuristic reduction of the search language.
It may be an excellent interactive choice, but it cannot justify a global
negative claim about all imported declarations. A resumable exhaustive mode
needs a fair route to the remaining permitted providers. An interactive receipt
should name the actual provider set searched.

\subsection{R7.2. Has type-guided abstraction refinement been combined with type classes?}
\answer Yes. This question has a direct answer in the proposal's own cited
Hoogle+ paper. Its Sections 2.5 and 5.1 describe type-class support by dictionary
passing: class constraints become extra dictionary arguments, class methods
consume dictionaries, and instances supply or construct them. The described
implementation does not cover arbitrary higher-kinded classes such as
\code{Monad}; its supported fragment should not be confused with full Lean type
classes \cite{hoogle}.

For Lean, represent an obligation such as \code{[C A]} as an explicit prerequisite
in the abstract application. A local instance is a currently available
resource. A global instance declaration is a dictionary-producing component,
possibly requiring other dictionaries. Concrete instance synthesis then checks
that the prerequisite can actually be discharged in the current environment.
A dictionary is not necessarily proof-irrelevant: classes can contain
computational operations. Erasing all class arguments would lose behavior as
well as applicability information.

The memoization key must include the class goal after relevant assignments,
local instances and their priorities, the environment, and instance-search
options. A cache indexed only by a printed pair ``class, type'' can be wrong in
a changed local context. A type metavariable may also make an instance goal
blocked rather than impossible; do not store that state as a permanent negative.
Universe metavariables, output parameters, and recursive instance prerequisites
are reasons to retain Lean's actual constraint machinery behind the abstract
index rather than recreating resolution in the net.

\recommend Adapt the existing dictionary-passing idea, then measure where Lean's
extra constraints cause false positives. This is an adaptation problem with a
clear prior solution for a smaller language, not a completely unexplored
combination of abstraction refinement and type classes.

\subsection{R7.3. Can premise selection use data goals and examples, and what corpus is suitable?}
\answer A data expression can certainly be supplied to a proposed ranking adapter,
but the inspected proof-premise systems do not establish that such a query plus
examples will produce good program-component rankings. LeanDojo and the Lean
hammer provide relevant evidence for proof-oriented retrieval, not a published
accuracy figure for this different task \cite{leandojo,leanhammer}.

Keep the interface small: a typed goal, local telescope, observation features,
provider metadata, and a fixed maximum number of suggestions. For proposition
goals, existing indexed lookup and proof-premise ranking are the natural first
baseline. For data goals, use structural type fit, reachability, selected
behavioral features, and only then learned ranking. A model trained to retrieve
proof lemmas may be useful, but it needs a separately measured data-goal track.

Body-erased definitions from Std, Lean core, and selected Mathlib modules are a
reasonable corpus source. Erasure must mean more than hiding the displayed body.
If the target declaration remains available as a provider, the answer is already
in the environment. Aliases, wrappers, an equation theorem that reconstructs
the desired body, or a helper copied from the target can leak the solution too.
The task harness should freeze an allowed provider set and audit its dependency
closure against the erased target and forbidden reconstructions.

A hidden reference implementation may serve as an external observation oracle
for the benchmark, but it must not appear in the synthesizer's term-building
context. Specifications should be published separately from that oracle. Record
which library lemmas are intentionally permitted; otherwise a task may change
silently from synthesis into lookup. Use separate tracks for provider-given,
retrieval-required, and helper-invention problems.

Split train and test data by function families and dependency clusters, not
only by declaration name. A test task that differs from a training task merely
by namespace or one wrapper is poor evidence of generalization. Report provider
recall at $k$, number of concrete applicability checks, dictionary-solving cost,
and final certified-program success. Good retrieval recall can coexist with
poor synthesis, and a fast synthesizer can mask weak retrieval on an overly
small provider corpus.

\section{R8. Ranking and equivalence}
\label{sec:r8}

\subsection{R8.1. Is there a defensible formal objective for ranking inhabitants?}
\answer Yes, as a declared preference model; no, as a uniquely correct ordering
determined by an under-specified type and contract. Minimum description length
is defensible because it makes a prior explicit. It does not remove the need
to choose that prior, nor does it prove agreement with an individual user's
intent.

For a fixed, well-defined grammar distribution, one possible objective is
\[
 J(p)=-\log P(p\mid T,C,\mathcal L),
\]
where $\mathcal L$ identifies the provider library and grammar version. The
probability model must actually assign probabilities to complete programs, or
be an explicitly stated subprobability model, rather than merely attaching
arbitrary weights and calling them probabilities. A code-length formulation
can instead use an explicit prefix encoding of syntax and component names.
Helper definitions must be charged; otherwise an arbitrarily complex solution
can be hidden behind a newly invented short name.

The practical objective should respect a hierarchy. First enforce the requested
type, contract, and axiom policy. Then distinguish executable from merely
logical results where execution matters. Within accepted candidates, report
program complexity, relevant library familiarity, estimated runtime, and proof
cost separately before imposing a configurable or project-fixed ordering.
A default order is a product decision; the report should not invent empirical
user-preference evidence for it.

Unused-input counts are at most a soft feature. A projection or constant
function can be exactly what the specification requests. Conversely, a term
can mention an input in $x-x$ without making its output depend on that input.
Syntactic occurrence is neither semantic information flow nor a correctness
criterion. Penalizing all ignored induction hypotheses likewise disfavors
legitimate predicates and short-circuiting traversals.

The present implementation ranks by unused binders, eliminator count, and raw
expression size, then deduplicates by pretty-printed form \cite{repo}. Replace
the last step independently of changing the preference model. Introduce an
explicit total tie-break on a stable structural representation; the objective
and tie-break belong in the receipt. Learning may change search priority while
$J$ remains frozen, as in R1. An optimality claim then refers to $J$ and the
covered grammar, not to an informal promise that the first program is what the
user intended.

\subsection{R8.2. Which equivalence is cheap and meaningful for dependent candidates?}
\answer Use a hierarchy of equivalence relations and label the level. Do not force
one relation to serve kernel identity, pruning, display, and user explanation.
The least expensive level is structural equality after safe normalization of
binding names and metadata, with all relevant universe and type information
retained. The next level is checked definitional equality or a certified
normalization. A stronger level is an explicit theorem of equality, possibly
pointwise equality of functions under the permitted axiom policy. Finite
observation agreement is a separate, weaker classification.

For two functions, equality on a probe set says only that the probes do not
separate them. It is useful for display: show one representative, the number of
others, and a distinguishing input whenever one is found. It is not a proof of
extensional equality. For dependent functions, probes and comparisons must be
well-typed in each result fiber; comparing values across fibers requires a
specified transport or heterogeneous relation.

Permanent observational pruning is justified only when the relation is a
congruence for every continuation that the remaining search can build, with
all future observations accounted for. This is a much stronger premise than
agreement on the current examples. If a future recursive call or synthesized
context can test a new input, preserve alternative representatives or refine
the equivalence classes. The identity-versus-length-cutoff example in R2
illustrates exactly how finite agreement can fail to extend.

Proof irrelevance helps erase distinctions between proofs of the same
proposition. It does not license dropping type dependencies, dictionaries that
carry computations, or transport endpoints. For program-size ranking it is
reasonable to count proof terms separately, but the checked term must keep the
evidence necessary for typing.

Finally, pretty printing is a presentation function, not an injective encoding
of terms. Hidden implicits, notation, and printing options can make distinct
terms look the same. Use a canonical expression serialization with structural
collision checks for exact deduplication. Use observation groups only for the
user-facing summary unless the stronger congruence theorem has been established.

\section{R9. Learned proposals and auditability}
\label{sec:r9}

\subsection{R9.1. What accuracy do models achieve on carrier and helper proposals?}
\answer The sources inspected for this report do not supply a reliable numerical
answer for the exact task: a Lean type and examples, a fixed carrier/helper
language, and an interactive end-to-end synthesis budget. Results for whole
proof generation or premise selection are not substitutes for that measurement.
The question is directly testable, and the benchmark should separate proposal
quality from the downstream search's ability to exploit a good proposal.

For each task, record an allowed carrier grammar, a known sufficient carrier or
helper specification, the permitted components, and a reference implementation
that is hidden from synthesis. Measure at least four successive events: the
proposal parses; it is well-typed and within the declared grammar; it enables
a solution under a fixed body-search work budget; and that solution passes the
full contract gate and execution checks. The first two are cheap interface
successes, not evidence of semantic success.

Compare four arms with the same total budget: the ordinary ordered grammar;
a fixed learned reordering of that grammar; model proposals restricted to the
grammar; and proposals allowed to expand it. The fourth arm changes the search
language and must be reported separately. Include a carrier-given oracle arm to
estimate how much headroom is actually attributable to carrier choice rather
than body search, proof automation, or unavailable components.

Report success at a fixed proposal count and at a fixed total cost. A top-ten
proposal list is not free if each item triggers an expensive search. Include
model inference time, rejected proposal processing, compilation and proof
costs, and time to first certified solution. Record exact model identifiers,
prompts, provider sets, seeds where available, and raw outputs. Use held-out
function families, several example sets per task, and uncertainty intervals
clustered by task family. Avoid treating many nearly identical variants as
independent observations.

The resulting study could justify a model lane without making it part of the
trusted logic. Until it is performed, the appropriate answer is an evaluation
protocol, not an invented percentage. This report's Python checks contain no
model-accuracy experiment and make no claim about one.

\subsection{R9.2. What determinism and auditability are possible beyond recording proposals?}
\answer Separate certificate replay, search replay, and model regeneration. A
closed accepted program and proof can be rechecked without reproducing any
search or model output. Replaying the same search additionally requires the
logical schedule and all proposal inputs. Regenerating identical model outputs
requires a controlled inference implementation and environment; a fixed seed
or temperature zero alone is not a sufficient specification of that system.

Use a restricted proposal protocol: a typed carrier expression, a list of local
variables to generalize, a helper type, or references to permitted component
names. Elaborate these in a frozen environment under explicit size and resource
limits. Do not accept an arbitrary tactic, import command, process invocation,
or metaprogram merely because it arrived in the same text field as a type.
The model chooses data for the synthesis engine; it does not choose the engine's
trust policy or acquire permission to execute arbitrary code.

A reproducible proposal record should contain the query and environment
fingerprints, grammar and provider versions, prompt bytes, model identifier,
inference settings, raw response bytes, parsed proposal representation,
validation outcome, and the logical point at which the proposal was admitted.
A stable canonical representation is useful for comparison, but preserve the
original bytes for auditing parser and elaborator behavior. Record absence or
cancellation of a proposal as an input event too.

If suggestions only reorder alternatives within a finite grammar and the
scheduler remains fair, the bounded completeness argument can be retained.
They can still change which candidate is discovered before an interactive
deadline. A model that proposes an out-of-grammar helper changes the grammar;
no receipt should describe that run as exhausting the original grammar under
an unchanged meaning of failure.

For a core-only default, the same proposal interface can be implemented by a
small deterministic grammar. An optional model is then a replaceable proposer,
not an architectural dependency. The acceptance receipt cites the checked
program, contract proof, and axiom inventory. Model output is useful provenance,
not a premise in the mathematical proof of correctness.

\section{An integrated architecture and implementation plan}
\label{sec:architecture}

\subsection{Six interfaces with explicit trust boundaries}
The recommendations above can be implemented without a single monolithic
``intelligent search'' subsystem. Six interfaces are enough to make the critical
contracts visible.

\paragraph{Frozen query.}
Contains the target and contract, local telescope or its closed abstraction,
imports and environment identity, axiom policy, provider set, grammar version,
ranking objective, and operational bounds. A changed provider set or equality
policy creates a new query identity even when the printed target is unchanged.

\paragraph{Persistent search node.}
Contains the relevant Lean saved state, ordered obligations, dependency summary,
consumed recursion resources, residual assumptions, ranking lower bound, and
stable decision-path identity. Global work accounting stays monotone outside
rollback; branch-local semantic data either rolls back or is guarded by complete
read-dependency keys.

\paragraph{Proposal service.}
Returns typed alternatives with provenance and a declared grammar effect.
Grammar enumeration, library retrieval, failed-proof generalization, and model
suggestions are implementations of this interface. The service may be wrong;
the consumer validates its output.

\paragraph{Observation service.}
Returns a typed observation result, dependencies, and either a conservative
refutation justification, a checked proof, or an unknown/residual constraint.
Contract consequences and speculative angelic assumptions have different
constructors. An evaluation error cannot masquerade as a Boolean false.

\paragraph{Proof service.}
Receives an obligation in a specific telescope and snapshot, a permitted axiom
profile, and a work quantum. It returns a checked proof, a validated
counterexample when appropriate, an unsupported-fragment report, or an
inconclusive resource result. It does not return a bare Boolean whose meaning
mixes these cases.

\paragraph{Acceptance and publication.}
The existing gate is the starting point for type, proof, and axiom checks.
Publication additionally verifies the executable representation when one is
promised. A receipt carries the exact certified term and enough provenance to
recheck it without trusting search, caches, models, or pretty printers.

\subsection{A proposed control loop}
The following is architectural pseudocode, not a compiled Lean implementation.
Its important feature is the placement of semantic decisions.
\begin{lstlisting}
freeze query, grammar, providers, trust policy, ranking objective
frontier := [initial well-scoped search node]
while frontier is not empty and logical budget remains:
    node := select by frozen epoch policy and stable tie-break
    restore node's saved Lean state
    validate or invalidate node's cached dependency records
    discharge eligible cheap constraints
    if there is a certified conservative refutation:
        discard this branch with its justification
    else if program and required proof are complete:
        run acceptance gate
        if accepted: record certificate; test publication separately
    else:
        choose an eligible obligation from the dependency graph
        request bounded proposals
        validate each proposal in a transaction
        enqueue resulting states, including the fair fallback
    commit learning and parallel results only at epoch boundaries
return accepted certificates plus exact coverage/stop metadata
\end{lstlisting}
Actual wall-clock cancellation can interrupt this loop or an individual service.
It must be reported separately from logical exhaustion. Long kernel or tactic
operations require their own interruption or heartbeat limits; checking the
clock only when entering a search node is not a hard bound on total runtime.
The current engine disables its outer heartbeat limit and uses explicit lane
deadlines, so this interaction deserves an integration test \cite{repo}.

\subsection{A revised sequence of milestones}
The proposal's research order is broadly sensible, but some low-risk reuse should
move earlier and some acceptance gates should change. In particular, receipts,
negative-result semantics, and rollback-safe cache identities are prerequisites
for trustworthy experimentation, not merely finishing features.

\begin{longtable}{@{}L{18mm}L{62mm}L{68mm}@{}}
\toprule
Phase & Implementation scope & Exit evidence\\
\midrule\endhead
P0 & Freeze source and benchmark manifests; define receipt, stop reasons,
cache ownership, and trust profiles & Recheck saved certificates without search;
distinguish unsupported, interrupted, grammar-exhausted, and proved-impossible\\
P1 & Separate logical acceptance from executable publication; indexed provider
lookup; guarded Nat/indexed induction; compare existing certified evaluation &
Compile/evaluate raw-recursion replacements; targeted probes; same bounded
semantic outcomes; measurements including failed attempts\\
P2 & First-class whole-state frontier, stable ordering, logical budgets,
dependency summaries, and conservative memoization & Exhaustive small-grammar
comparison against the reference search; replay identical decision traces;
rollback adversarial tests\\
P3 & Typed per-observation reports, certified decompositions, guarded incremental
caches, local proof obligations & Cache-on/cache-off semantic equivalence;
all fragment refutations justified; local induction-hypothesis tests\\
P4 & Bounded carrier packages, conflict-graph lower bounds, recursion capabilities,
helper/generalization proposals & Carrier-given versus invented tracks;
finite completeness checks; interpretation/realization proof obligations\\
P5 & Richer motives and transport policies, inductive contracts, refined
library reachability & Indexed and quantified-contract coverage; provider recall;
explicit limits of each supported fragment\\
Optional & Parallel speculative work and learned proposal sources behind the same
interfaces & Deterministic epoch replay; end-to-end latency/cost ablations;
no change in certificate acceptance semantics\\
\bottomrule
\end{longtable}
Nat recursion and basic guards need not wait for an invented minimal-carrier
theory. Equally, a custom evaluator should not be justified merely by the fact
that residual evaluation was expensive on one historical query. Measure the
existing proof-producing evaluator and finer observation decomposition first.
Retrieval of a missing standard lemma can be independent of the carrier work.

The acceptance criterion ``identical candidate lists under equal wall time''
should be removed from the search-rewrite phase. It is neither a meaningful
semantic equivalence test nor stable across machines. Instead compare complete
candidate sets on a finite grammar, first results under a fixed logical trace,
and independently measured latency distributions. Ranking changes should have
a separate benchmark and explicit version.

\subsection{Engineering items that still need semantic care}
The engineering/research distinction is useful, but the proposal's engineering
items are not exempt from correctness obligations. E1 must link a compilable
presentation to the certified program. E2 must preserve local scoping and
expected-type constraints when turning a sketch into goals. E3 must record the
actual environment and choices, not just a printed answer. E4 needs a profile
that attributes work without counting nested timers twice. E5 needs deterministic
commit rules, rather than a race followed by sorting. E7 must preserve trust
profiles when adding tactics, and E8 must prevent benchmark leakage.

E6 has an additional abstraction issue. A propositional decision procedure over
independent atoms can produce an intuitionistic proof or a countermodel for the
\emph{reified formula}. That countermodel need not refute the original Lean goal
when the chosen atoms are concrete propositions related by arithmetic or
contextual hypotheses. For example, treating a theorem of arithmetic as an
unconstrained atom can make an unprovable propositional abstraction of a
provable Lean statement. Thus report ``not derivable in this propositional
abstraction,'' unless an interpretation theorem transfers the countermodel to
the actual query. A kernel-checked proof of $T\to\mathrm{False}$ remains a
fundamentally different artifact.

Finally, ``no settings'' for an ordinary user does not mean ``no configuration
identity.'' Defaults, imported attributes, grammar bounds, axiom profiles, and
model versions all affect behavior. Keep the interactive surface simple while
making those choices explicit in receipts and benchmark manifests.

\section{Validation, experiments, and limits of the evidence}
\label{sec:validation}

\subsection{What was executed for this report}
The accompanying \path{experiments/model_checks.py} is a deterministic,
standard-library-only Python program. All eight groups of assertions passed.
It explores deliberately small models of the semantic distinctions above; it
is neither a port of \repo{} nor a Lean kernel verification.

\begin{longtable}{@{}L{32mm}L{66mm}L{50mm}@{}}
\toprule
Check & Executed result & What it establishes\\
\midrule\endhead
Coupled-goal cost & Two obligations close in one action of cost 1; summing unit
bounds gives 2 & Counterexample to the unqualified open-goal sum heuristic\\
Reverse fold & All 255 binary lists of length at most 7 agree with reverse &
Finite validation of the list-carrier implementation, not a complexity claim\\
Incomplete rows & Three distinct partial rows have only one conflict edge;
chromatic number 2 & Distinct partial rows cannot be counted as distinct states\\
Finite carriers & Exhaustive parity and unary length-modulo-three tables;
minimum sample-consistent state counts 2 and 3 & Exact minima in the declared
finite table models\\
Rollback cache & Hole-only cache returns stale value 0; branch-versioned cache
returns expected value 1 & A minimal example of the need for assignment identity\\
Parametric probes & Two uniform list functions agree at lengths 0 through 4 and
differ at length 5 & Finite length coverage does not prove universal equality\\
Deadline versus work & Simulated work limits 10 and 20 discover different sets;
fixed work 10 gives the same set & Sorting cannot repair different discovery sets\\
Conjunctive retrieval & An ordinary path reaches $C$ from $A$ for a provider
requiring $A$ and $B$; the hypergraph does not & Application prerequisites must
be conjunctive\\
\bottomrule
\end{longtable}
The finite-carrier enumeration is small enough to inspect exactly. It fixes the
alphabet and output sets and enumerates all labeled transition tables, all
seeds, and all finalizers at each stated state count.

\begin{center}
\begin{tabular}{@{}lrrrr@{}}
\toprule
Sample task & Samples & States & Tables checked & Solutions\\
\midrule
Binary parity, length $\leq4$ & 31 & 1 & 2 & 0\\
 & 31 & 2 & 128 & 2\\
Unary length modulo 3, length $\leq5$ & 6 & 1 & 3 & 0\\
 & 6 & 2 & 72 & 0\\
 & 6 & 3 & 2187 & 6\\
\bottomrule
\end{tabular}
\end{center}
The 2392 table checks confirm the executable finite procedure in R2. The first
parity witness uses seed $0$, finalizer $(0,1)$, and transitions that leave the
state unchanged on symbol $0$ and exchange it on symbol $1$. The first unary
witness cycles $0\mapsto1\mapsto2\mapsto0$ and returns the current state.
The script's solution counts include different labelings; it does not quotient
machines by isomorphism.

The rollback check is a model counterexample, not a reproduced defect in the
actual Lean implementation. Likewise, the deadline simulation uses logical
work positions, not measured processors or seconds. These distinctions are
part of the results file itself, not qualifications added only to this article.
Run the checks from the archive root with:
\begin{lstlisting}
python3 experiments/model_checks.py --output experiments/results.json
\end{lstlisting}

\subsection{The required Lean regression and ablation suite}
A production implementation should add tests at three different levels.
First, semantic unit tests validate individual rules: scope preservation,
metavariable isolation, context substitution, contract decomposition,
conservative false results, transport reconstruction, and recursive-call
capabilities. Every cache should have a cache-disabled reference mode.

Second, exhaustive bounded tests compare the original and refactored search on
small finite provider and syntax sets. Compare typed candidate sets before
ranking. Exercise shared universe variables, delayed assignments, changed local
instances, identical-looking goals in different telescopes, and rollback to a
branch that reuses an identifier with a different assignment. The acceptance
gate does not detect a valid candidate that an unsound prune silently discarded,
so positive-result testing alone is inadequate.

Third, end-to-end benchmarks measure usefulness. Preserve the existing baseline
as historical regression material, but add indexed families, Nat recursion,
conditional programs, carrier-given and carrier-invented Church problems,
quantified contracts, and retrieval-required definitions. Each entry records
its source, permitted providers, examples, universal contract where available,
reference proof or implementation, and leakage controls.

Use both a logical-work track and a wall-clock track. The first tests replay and
coverage; the second tests responsiveness at declared budgets such as 1, 5, 10,
and 60 seconds. Those budgets are proposed measurement points, not results of
this report. Distinguish cold environment/index loading from a warmed session.
Report time to the first certified candidate, number and rank of distinct
behavioral groups, proof and compilation cost, peak memory, and unresolved
versus refuted outcomes.

An ablation should change only one mechanism at a time: fine-grained
observations, existing \code{cbv}, guarded caching, dependency-aware goal order,
carrier proposals, retrieval, or learned priorities. The proposal's historical
45\% residual-evaluation figure is motivation to measure evaluation, not proof
that replacing that evaluator improves all queries \cite{repo}.

\subsection{Claims that this investigation does not make}
No Lean executable was available in the execution environment, so no revised
Lean engine, tactic wrapper, or example program was compiled here. The pinned
repository and selected Lean source files were inspected through their source
interfaces. The repository's reported benchmark scores remain its historical
measurements, not independently reproduced results.

The propositions and theorem in this report are mathematical arguments with
explicit hypotheses, not machine-checked Lean theorems. Their small-model
instances were tested where stated. In particular, the dependency analysis and
interpreter simulation still require implementation-specific proofs or checked
certificates. No new asymptotic complexity bound for unrestricted dependent
synthesis, no universal minimal-carrier algorithm, and no model-accuracy number
is claimed.

Bibliographic searches establish positive evidence for the cited mechanisms.
They do not establish that no other paper has addressed an adjacent problem.
Statements that a source does not answer a question refer to the inspected
sources and the precisely stated task. This is especially important for
higher-order dependent live evaluation, automated generalization, and
latency-bounded model proposals, where neighboring research can use a different
language or contract model.

\section{Conclusion}
\label{sec:conclusion}
The central expert answer is that \repo{} does not need one breakthrough that
solves all nine research areas simultaneously. It needs explicit boundaries
between an accepted program, a conservative search restriction, a bounded
negative result, and a preference among solutions. Once those boundaries are
fixed, several apparently open implementation questions have immediate starting
points in Lean and the existing synthesis literature.

Whole-state snapshots provide a sound baseline for a dependency-aware search;
reproducible logical epochs provide a meaningful alternative to impossible
machine-independent deadline promises. Finite carrier packages can be searched
exactly, and witnessed incompatibility gives a sound lower bound without
pretending that incomplete observation rows form a known minimal automaton.
Typed hole closures, guarded incremental caches, and equation-based certified
evaluation support stronger pruning without treating a final kernel gate as a
repair for lost solutions. Motive computation, indexed induction, and dictionary
constraints have existing machinery to reuse. Program-and-proof synthesis can
start from local verification conditions instead of waiting for universal lemma
invention.

The recommended next artifact is therefore an auditable bounded engine with
strong small-instance tests, not an unqualified ``general synthesis'' claim.
Its research extensions can then be evaluated independently: richer carrier
proposals, better generalization, dependent residual constraints, and optional
learned guidance. Each improvement has a concrete success criterion, a known
failure mode, and a receipt that states exactly what was checked.

\appendix
\section{Source map and question coverage}
The numbering R1.1--R9.2 in this report follows the order of the questions in the
proposal, not a new selection of topics. The following map makes it possible to
compare the answers directly with the original source at the pinned commit.
Paths in the first nine rows are relative to
\path{docs/proposals/11-further-improvements/sections/}.

\begin{longtable}{@{}L{39mm}L{64mm}L{45mm}@{}}
\toprule
Source & Questions covered & Report section\\
\midrule\endhead
\path{05-search.tex} & R1.1--R1.4: independence, lower bounds, online costs,
wall-clock determinism & \ref{sec:r1}, p.~\pageref{sec:r1}\\
\path{06-carriers.tex} & R2.1--R2.5: minimal carriers, polymorphism, fusion,
recursor motives, generalization & \ref{sec:r2}, p.~\pageref{sec:r2}\\
\path{07-evaluation.tex} & R3.1--R3.4: higher-order/dependent holes,
incrementality, verified evaluation, angelic completeness & \ref{sec:r3},
p.~\pageref{sec:r3}\\
\path{08-dependent.tex} & R4.1--R4.4: motive enumeration, transport,
transparency, Lean API & \ref{sec:r4}, p.~\pageref{sec:r4}\\
\path{09-recursion.tex} & R5.1--R5.3: realization, top-down angelic execution,
scheme basis & \ref{sec:r5}, p.~\pageref{sec:r5}\\
\path{10-contracts.tex} & R6.1--R6.4: proof/value scheduling, inductive fragment,
Lean automation, proof debt & \ref{sec:r6}, p.~\pageref{sec:r6}\\
\path{11-retrieval.tex} & R7.1--R7.3: dependent abstraction, type classes,
program-component relevance & \ref{sec:r7}, p.~\pageref{sec:r7}\\
\path{12-ranking.tex} & R8.1--R8.2: ranking objectives, equivalence &
\ref{sec:r8}, p.~\pageref{sec:r8}\\
\path{12-ranking.tex} & R9.1--R9.2: model proposal accuracy, auditability &
\ref{sec:r9}, p.~\pageref{sec:r9}\\
\midrule
\path{Leant2/Search/Core.lean} & Search rules, residual evaluation, induction
filter, proof portfolio & Sections \ref{sec:r1}--\ref{sec:r7}\\
\path{Leant2/Native/Transaction.lean} & Saved-state boundary, global ledger,
blocker cache, deadlines & Sections \ref{sec:r1}, \ref{sec:r3}\\
\path{Leant2/Engine.lean} & Candidate ranking, grace period, contract extraction &
Sections \ref{sec:r1}, \ref{sec:r8}\\
\path{Leant2/Accept/Gate.lean} & Kernel checking and transitive axiom audit &
Sections \ref{sec:foundations}, \ref{sec:architecture}\\
\path{Leant2/Core/Types.lean} & Trust profiles and result algebra &
Sections \ref{sec:foundations}, \ref{sec:architecture}\\
\bottomrule
\end{longtable}
The proposal's evidence, engineering, adaptation, roadmap, expert-summary, and
bibliography sections were also reviewed. In particular, this report does not treat the proposal's
explicitly memory-based bibliography as already verified: the relevant positive
claims were checked against the primary sources listed below.

\clearpage
\phantomsection
\addcontentsline{toc}{section}{References}
\begin{thebibliography}{99}
\small
\bibitem{repo}
Vladimir Reshetnikov. \emph{Leant2 repository}, snapshot
\code{784664ff34e819422510a4d470ab07fe48bb736b}, 18 September 2026.
\emph{Leant2: Design for Further Improvements} and implementation files mapped
in Appendix A. Accessed 21 September 2026.
\url{https://github.com/VladimirReshetnikov/Leant2/tree/784664ff34e819422510a4d470ab07fe48bb736b}.

\bibitem{leaninduction}
Lean contributors. \emph{Induction metaprogramming implementation}, Lean 4.34.0.
Source-level API inspection, accessed 21 September 2026.
\url{https://github.com/leanprover/lean4/blob/v4.34.0/src/Lean/Meta/Tactic/Induction.lean}.

\bibitem{leancbv}
Lean contributors. \emph{Call-by-value tactic frontend}, Lean 4.34.0.
Includes \code{cbv} and \code{decide_cbv}. Accessed 21 September 2026.
\url{https://github.com/leanprover/lean4/blob/v4.34.0/src/Lean/Elab/Tactic/Cbv.lean}.

\bibitem{leantactics}
Lean contributors. \emph{The Lean Language Reference: Tactic Reference}.
Moving online reference, accessed 21 September 2026. Capability descriptions
from the current reference are distinguished from pinned source checks.
\url{https://lean-lang.org/doc/reference/latest/Tactic-Proofs/Tactic-Reference/}.

\bibitem{leanrec}
Lean contributors. \emph{The Lean Language Reference: Recursive Definitions}.
Accessed 21 September 2026.
\url{https://lean-lang.org/doc/reference/latest/Definitions/Recursive-Definitions/}.

\bibitem{leaninductbook}
Lean contributors. \emph{Theorem Proving in Lean 4: Inductive Types}.
Natural-number recursion and addition. Accessed 21 September 2026.
\url{https://lean-lang.org/theorem_proving_in_lean4/Inductive-Types/}.

\bibitem{aesop}
Jannis Limperg and Asta Halkj\ae r From. Aesop: White-Box Best-First Proof Search
for Lean. \emph{CPP}, 2023.
\url{https://doi.org/10.1145/3573105.3575671}.
Author-deposited record: \url{https://zenodo.org/records/7430233}.

\bibitem{aesopreadme}
Aesop contributors. \emph{Aesop README}, accessed 21 September 2026.
Source blob \path{f75ec1f8bebcf6d857fc3954cf91cf99f323a758}.
\url{https://github.com/leanprover-community/aesop/blob/master/README.md}.

\bibitem{canonical}
Chase Norman and Jeremy Avigad. Canonical for Automated Theorem Proving in Lean.
\emph{ITP}, 2025. Inspected preprint version 1.
\url{https://arxiv.org/html/2504.06239v1}.

\bibitem{probe}
Shraddha Barke, Hila Peleg, and Nadia Polikarpova. Just-in-Time Learning for
Bottom-Up Enumerative Synthesis. \emph{OOPSLA}, 2020.
\url{https://arxiv.org/abs/2010.08663}.

\bibitem{luby}
Michael Luby, Alistair Sinclair, and David Zuckerman. Optimal Speedup of Las Vegas
Algorithms. \emph{Information Processing Letters} 47(4), 173--180, 1993.
\url{https://www.sciencedirect.com/science/article/pii/0020019093900299}.

\bibitem{gibbons}
Jeremy Gibbons, Graham Hutton, and Thorsten Altenkirch. When is a Function a Fold
or an Unfold? \emph{Electronic Notes in Theoretical Computer Science} 44(1),
146--160, 2001.
\url{https://www.cs.ox.ac.uk/publications/publication2368-abstract.html}.

\bibitem{weber}
Tjark Weber and James Caldwell. Constructively Characterizing Fold and Unfold.
\emph{LOPSTR 2003}, LNCS 3018, 110--127.
\url{https://doi.org/10.1007/978-3-540-25938-1_11}.

\bibitem{polytesting}
Jean-Philippe Bernardy, Patrik Jansson, and Koen Claessen. Testing Polymorphic
Properties. \emph{ESOP}, LNCS 6012, 125--144, 2010.
\url{https://doi.org/10.1007/978-3-642-11957-6_8}.

\bibitem{angluin}
Dana Angluin. Learning Regular Sets from Queries and Counterexamples.
\emph{Information and Computation} 75(2), 87--106, 1987.
\url{https://doi.org/10.1016/0890-5401(87)90052-6}.

\bibitem{smyth}
Justin Lubin, Nick Collins, Cyrus Omar, and Ravi Chugh. Program Sketching with
Live Bidirectional Evaluation. \emph{Proceedings of the ACM on Programming
Languages} 4(ICFP), article 109, 2020.
\url{https://doi.org/10.1145/3408991}.
Preprint: \url{https://arxiv.org/abs/1911.00583}.

\bibitem{scrybe}
Niek Mulleners, Johan Jeuring, and Bastiaan Heeren. Program Synthesis Using
Example Propagation. Preprint, 2022.
\url{https://arxiv.org/abs/2210.13873}.

\bibitem{selfadjust}
Umut A. Acar, Matthias Blume, and Jacob Donham. A Consistent Semantics of
Self-Adjusting Computation. Preprint, 2011.
\url{https://arxiv.org/abs/1106.0478}.

\bibitem{lean4lean}
Mario Carneiro. Lean4Lean: Verifying a Typechecker for Lean, in Lean.
Preprint, version 3, revised 14 September 2025.
\url{https://arxiv.org/abs/2403.14064v3}.

\bibitem{burst}
Anders Miltner, Adrian Trejo Nu\~nez, Ana Brendel, Swarat Chaudhuri, and Isil Dillig.
Bottom-Up Synthesis of Recursive Functional Programs using Angelic Execution.
\emph{Proceedings of the ACM on Programming Languages} 6(POPL), article 21, 2022.
\url{https://doi.org/10.1145/3498682}.
Preprint: \url{https://arxiv.org/abs/2107.06253}.

\bibitem{sauto}
\L ukasz Czajka. Practical Proof Search for Coq by Type Inhabitation.
\emph{IJCAR}, 28--57, 2020.
\url{https://doi.org/10.1007/978-3-030-51054-1_3}.
Author page: \url{https://www.mimuw.edu.pl/~lukaszcz/sauto/index.html}.
Unfolding options checked in the project documentation:
\url{https://coqhammer.github.io/}.

\bibitem{synquid}
Nadia Polikarpova, Ivan Kuraj, and Armando Solar-Lezama. Program Synthesis from
Polymorphic Refinement Types. \emph{PLDI}, 2016.
\url{https://arxiv.org/abs/1510.08419v3}.

\bibitem{hoogle}
Zheng Guo, Michael James, David Justo, Jiaxiao Zhou, Ziteng Wang, Ranjit Jhala,
and Nadia Polikarpova. Program Synthesis by Type-Guided Abstraction Refinement.
\emph{Proceedings of the ACM on Programming Languages} 4(POPL), article 12, 2020.
\url{https://doi.org/10.1145/3371080}.
Inspected author PDF: \url{https://ranjitjhala.github.io/static/hoogle_plus.pdf}.

\bibitem{trio}
Woosuk Lee and Hangyeol Cho. Inductive Synthesis of Structurally Recursive
Functional Programs from Non-recursive Expressions.
\emph{Proceedings of the ACM on Programming Languages} 7(POPL), article 70, 2023.
\url{https://doi.org/10.1145/3571263}.

\bibitem{leandojo}
Kaiyu Yang et al. LeanDojo: Theorem Proving with Retrieval-Augmented Language
Models. \emph{NeurIPS}, 2023.
\url{https://arxiv.org/abs/2306.15626}.

\bibitem{leanhammer}
Thomas Zhu, Joshua Clune, Jeremy Avigad, Albert Q. Jiang, and Sean Welleck.
Premise Selection for a Lean Hammer. Preprint, 2025.
\url{https://arxiv.org/abs/2506.07477}.
Author project page: \url{https://cmu-l3.github.io/lean-hammer/}.
\end{thebibliography}
\end{document}
